\subsection{维数与截切序列}

设 A 为诺特环，M 为有限生成 A-模，S 为 A 的根基的一个子集，S 为 A 中由 S 生成的理想，而 a 为 M 的湮灭子。我们有

\[
\dim_ {\mathrm{A}} (\mathrm{M}) = \dim (\mathrm{A} / \mathrm{a}) \leqslant \dim (\mathrm{A});\tag{6}
\]

因此，若 A 是局部环，则有 \(\dim_{\mathbf{A}}(\mathbf{M}) < +\infty\)。此外，根据 II, § 4, no. 4 中 Prop. 18 的推论，A-模 M/SM 的支集等于 V(a + S)，因而

\[
\dim_ {A} (M / S M) = \dim (A / (a + \mathfrak {S}))  .\tag{7}
\]

当 S 有限时，或者 M 不为 0 时，我们有不等式

\[
\dim_ {\mathbf {A}} (\mathbf {M} / \mathrm{SM}) \leqslant \dim_ {\mathbf {A}} (\mathbf {M}) \leqslant \operatorname{Card} (\mathbf {S}) + \dim_ {\mathbf {A}} (\mathbf {M} / \mathrm{SM});\tag{8}
\]

当 S 有限时，这由 no. 1 的 prop. 2 以及上面的公式 (6) 和 (7) 得出，而 S 无限的情形是平凡的。

定义 1. --- 设 A 为诺特局部环，M 为非零有限生成 A-模，S 为 A 的极大理想 \(\mathfrak{m}_{\mathrm{A}}\) 的子集。若有

\[
\dim_ {A} (M) = \operatorname{Card} (S) + \dim_ {A} (M / S M).\tag{9}
\]

若 S 是 M 的截切集，则有 \(\operatorname{Card}(S) \leqslant \dim_{\mathbf{A}}(\mathbf{M})\)，所以 S 是有限集。若族 \((x_i)_{i \in I}\) 中的元素属于 \(\mathfrak{m}_{\mathbf{A}}\)，且有

\[
\dim_ {\mathrm{A}} (\mathrm{M}) = \operatorname{Card} (\mathrm{I}) + \dim_ {\mathrm{A}} \left(\mathrm{M} / \sum_ {i \in \mathrm{I}} x _ {i} \mathrm{M}\right),\tag{10}
\]

也就是说，若 \(i\mapsto x_{i}\) 是从 I 到 \(\mathfrak{m}_{\mathbf{A}}\) 中某个对于 M 为截切的子集的双射。

若 \(x\) 是 \(m_A\) 中的元素，称其为 \(M\) 的截切元，当且仅当 \(\{x\}\) 是 \(M\) 的截切子集，即有

\[
\dim_ {\mathbf {A}} (\mathbf {M} / x \mathbf {M}) = \dim_ {\mathbf {A}} (\mathbf {M}) - 1.
\]

备注。--- 1) 由公式 (6) 和 (7) 可知，S 是 M 的截切集，当且仅当它是 A/a 的截切集，其中 a 是 M 的湮灭子。

\begin{enumerate}
\def\labelenumi{\arabic{enumi})}
\setcounter{enumi}{1}
\tightlist
\item
  设 S 和 S′ 是 \(m_{A}\) 的两个不相交子集。要使 \(S \cup S'\) 成为 M 的截切集，必要且充分的条件是 S 为 M 的截切集，且 S' 为 \(M' = M/SM\) 的截切集。这由不等式 (8) 和公式
\end{enumerate}

\[
\begin{array}{r l} \operatorname{Card} (S \cup S ^ {\prime}) + \dim_ {A} (M / (S M + S ^ {\prime} M)) - \dim_ {A} (M) = & \\ = (\operatorname{Card} (S) + \dim_ {A} (M / S M) - \dim_ {A} (M)) + & \\ + (\operatorname{Card} (S ^ {\prime}) + \dim_ {A} (M ^ {\prime} / S ^ {\prime} M ^ {\prime}) - \dim_ {A} (M ^ {\prime})) \end{array}
\]

\begin{enumerate}
\def\labelenumi{\arabic{enumi})}
\setcounter{enumi}{2}
\tightlist
\item
  设 \(x \in \mathfrak{m}_{\mathbf{A}}\)，且 \(n \geqslant 1\) 为整数。显然，模 \(\mathbf{M} / x\mathbf{M}\) 和 \(\mathbf{M} / x^n\mathbf{M}\) 具有相同的支集，因而具有相同的维数。因此，\(x\) 是 \(\mathbf{M}\) 的截切元，当且仅当 \(x^n\) 是 \(\mathbf{M}\) 的截切元。由此以及备注 2，立即得到如下结果：设 \(x_1, ..., x_r\) 是 \(\mathfrak{m}_{\mathbf{A}}\) 中的元素，且 \(n_1, ..., n_r\) 是 \(> 0\) 的整数；则序列 \((x_1, ..., x_r)\) 是 \(\mathbf{M}\) 的截切序列，当且仅当序列 \((x_1^{n_1}, ..., x_r^{n_r})\) 是 \(\mathbf{M}\) 的截切序列。
\end{enumerate}

命题 3. --- 设 A 为诺特局部环，M 为非零有限生成 A-模。对于 \(\mathfrak{m}_{\mathrm{A}}\) 中的元素 x，x 是 M 的截切元的充要条件是：x 不属于任何极小元 \(\mathfrak{p}\)，其属于 \(\operatorname{Supp}(M)\) 且满足 \(\dim(\mathrm{A}/\mathfrak{p}) = \dim_{\mathrm{A}}(\mathrm{M})\)；而 M 上比率为 x 的倍乘映射 \(x_{\mathrm{M}}\) 为单射则是充分条件。

设 a 是 M 的湮灭子。说 x 是 M 的截切元，意味着 x 是 A/a 的截切元，而 M 的支集由 A 中满足 \(a \subset p\) 的素理想 p 构成；若 \(x_{M}\) 是单射，则 x 在 A/a 中的像不是 A/a 中的零因子。将 No. 1 的 Prop. 2 的 Cor. 2 应用于环 A/a，即得 Prop. 3。

推论。--- \(m_{A}\) 中对于 M 完全截切的每个元素序列（A, X, p. 157, def. 2）都是 M 的截切序列。

设 \((x_{1},\ldots ,x_{r})\) 是 \(\mathfrak{m}_{\mathbf{A}}\) 中对于 M 完全截切的元素序列。置 \(\mathbf{M}_0 = \mathbf{M}\)，并归纳定义 \(\mathbf{M}_i = \mathbf{M}_{i - 1} / x_i\mathbf{M}_{i - 1}\)，其中 \(1\leqslant i\leqslant r\)。根据 A, X, p. 160 的 cor. 1，比率为 \(x_{i}\) 且作用于 \(\mathbf{M}_{i - 1}\) 的倍乘映射是单射，因而 \(\dim_{\mathbf{A}}(\mathbf{M}_i) = \dim_{\mathbf{A}}(\mathbf{M}_{i - 1}) - 1\)（\(1\leqslant i\leqslant r\)）（prop. 3）。因此有

\[
\dim_ {\mathbf {A}} (\mathbf {M}) = r + \dim_ {\mathbf {A}} \left(\mathrm{M} / (x _ {1} \mathrm{M} + \dots + x _ {r} \mathrm{M})\right),
\]

所以序列 \((x_{1},...,x_{r})\) 是 M 的截切序列。

备注 4. --- 一般而言，M 的截切序列不一定是 M 的完全截切序列（p. 87, exerc. 6）。稍后我们将研究这样的诺特局部环上的模：其中每个截切序列都是完全截切的。

定理 1. --- 设 A 为诺特局部环，M 为非零有限生成 A-模，S 为 A 的极大理想 \(\mathfrak{m}_{\mathbf{A}}\) 的子集。

\begin{enumerate}
\def\labelenumi{\alph{enumi})}
\item
  若 M/SM 具有有限长度，则有 Card(S) \(\geqslant\) dim \(_{A}\) (M)；若 S 是 M 的截切集，则 Card(S) \(\leqslant\) dim \(_{A}\) (M)。
\item
  M 的每个截切子集都包含于 M 的某个极大截切子集中。
\item
  下列性质等价：
\end{enumerate}

\begin{enumerate}
\def\labelenumi{(\roman{enumi})}
\item
  S 是 M 的极大截切子集；
\item
  S 是 M 的截切子集，且 \(\operatorname{Card}(\mathbf{S}) = \dim_{\mathbf{A}}(\mathbf{M})\)；
\item
  M/SM 具有有限长度，且 Card(S) = dim\_A(M)；
\item
  S 是 M 的截切子集，且 M/SM 具有有限长度。
\end{enumerate}

由于 \(S \subset m_{A}\)，Nakayama 引理表明 M/SM ≠ \{0\}，因而 dim \(_{A}\) (M/SM) ≥ 0，且等号成立当且仅当 M/SM 具有有限长度。断言 \(a\) ) 随即由公式 (8)、(9) 以及性质 (ii)、(iii)、(iv) 的等价性得到。

断言 \(b\) ) 由如下事实得到：\(\mathfrak{m}_{\mathbf{A}}\) 中每个对于 \(\mathbf{M}\) 为截切的子集的基数都不超过整数 \(\dim_{\mathbf{A}}(\mathbf{M})\)。

由 \(a\) )，每个基数等于 \(\dim_{\mathbf{A}}(\mathbf{M})\) 的 M 的截切子集都是极大的。还需证明：若 S 是 M 的截切集且 \(\operatorname{Card}(\mathbf{S}) < \dim_{\mathbf{A}}(\mathbf{M})\)，则 S 不是极大的。设 a 是 M 的湮灭子，B 是局部诺特环 A/(a + SA)。根据 n° 1 的 prop. 2 的 cor. 2，存在元素 \(x\)，属于 \(\mathfrak{m}_{A}\)，使得 \(\dim(\mathrm{B}/x\mathrm{B}) = \dim(\mathrm{B}) - 1\)；因而 \(x \notin \mathrm{S}\)。根据备注 2，子集 S ∪ \(\{x\}\) 的 \(\mathfrak{m}_{A}\) 是 A/a 的截切集，由备注 1，它也是 M 的截切集。

推论。--- M 的维数是所有满足下述条件的整数 \(d \geqslant 0\) 中的最小者：存在一个序列 \((x_{1}, ..., x_{d})\)，其元素属于 \(m_{A}\)，使 A-模 M/ \(\sum_{i=1}^{d} x_{i}M\) 具有有限长度。

由于 \(\varnothing\) 是 M 的截切子集，定理 1 的 \(b\) ) 证明了 M 存在极大截切序列，即 \((x_{1},\ldots,x_{d})\)。有 \(d=\dim_{\mathbf{A}}(\mathbf{M})\) 且 A-模 \(\mathbf{M}/\sum_{i=1}^{d}x_{i}\mathbf{M}\) 具有有限长度，根据定理 1 的性质 (iii) 和 \(c\) )。反之，若 \((x_{1}',\ldots,x_{d'}')\) 是 \(\mathfrak{m}_{\mathbf{A}}\) 中元素组成的序列，使 A-模 \(\mathbf{M}/\sum_{j=1}^{d'}x_{j}'\mathbf{M}\) 具有有限长度，则有 \(d'\geqslant\dim_{\mathbf{A}}(\mathbf{M})\)，由 th. 1 的 \(a\) )。

回忆（III, § 3, n° 2, def. 1）：诺特局部环 A 的理想 q 称为 A 的定义理想，如果 A 的 q-进拓扑与 m\_A-进拓扑重合。

引理 2. --- 设 A 为诺特局部环，q 为 A 的理想。下列条件等价：

\begin{enumerate}
\def\labelenumi{(\roman{enumi})}
\item
  q 是 A 的定义理想；
\item
  存在整数 \(n \geqslant 0\)，使得 \(\mathfrak{m}_{\mathbf{A}}^{n} \subset \mathfrak{q} \subset \mathfrak{m}_{\mathbf{A}}\)；
\item
  q ≠ A，且 A/q 是具有有限长度的 A-模；
\item
  V(q) 等于 \(\{m_{A}\}\)（换言之，\(m_{A}\) 是包含 q 的唯一 A 的素理想）。
\end{enumerate}

事实上，(i)、(ii) 和 (iv) 的等价性已在 III, § 2, no. 5 中证明，而 (i) 与 (iii) 的等价性由 IV, § 2, no. 5, cor. 2 to prop. 9 得出。

因此，Thm. 1 的推论使我们可以陈述如下附论：

附论。--- 诺特局部环 A 的维数是满足存在一个由 d 个元素生成的 A 的定义理想的整数 \(d \geqslant 0\) 中的最小者。

\subsection{首批应用}

设 A 为诺特环，V = Spec(A) 为其谱。在本节中，我们把 V 中任何形如 \footnote{回忆，\(V(x)\) 由 \(A\) 的包含 \(x\) 的素理想组成。} V(x) 且 \(x \in A\) 的子集称为超曲面。

命题 4. --- 设 X 是 V 的闭子集，\(H_{1}, \ldots, H_{m}\) 是 V 中的超曲面。置 \(X' = X \cap H_{1} \cap \ldots \cap H_{m}\)。

\begin{enumerate}
\def\labelenumi{\alph{enumi})}
\tightlist
\item
  对于包含于 X′ 的 V 的每个闭子集 Y，有
\end{enumerate}

\[
\operatorname{codim} (\mathrm{Y}, \mathrm{X} ^ {\prime}) \geqslant \operatorname{codim} (\mathrm{Y}, \mathrm{X}) - m.
\]

\begin{enumerate}
\def\labelenumi{\alph{enumi})}
\setcounter{enumi}{1}
\item
  对于 X' 的每个不可约分支 Z，有 codim(Z, X) ≤ m。若 X' 非空，则 codim(X', X) ≤ m。
\item
  若 Z 是包含于 X 的 V 的闭不可约子集，且 \(\text{codim}(Z, X) \leqslant m\)，则存在超曲面 \(H_{1}^{\prime}, \ldots, H_{m}^{\prime}\)，使 Z 成为 \(X \cap H_{1}^{\prime} \cap \ldots \cap H_{m}^{\prime}\) 的不可约分支。
\end{enumerate}

设 \(\mathfrak{a}\) 是 A 的理想，\(x_{1},\ldots,x_{m}\) 是 A 中的元素，使得 \(\mathrm{X}=\mathrm{V}(\mathfrak{a})\)，且 \(\mathrm{H}_{i}=\mathrm{V}(x_{i})\)（\(1\leqslant i\leqslant m\)）。设 Z 是包含于 X 的 V 的闭不可约子集；则存在包含 \(\mathfrak{a}\) 的 A 的素理想 p，使 \(\mathrm{Z}=\mathrm{V}(\mathfrak{p})\)。

首先假设 Z 包含于 X'，并记 \(\xi_{i}\) 为 \(x_{i}\) 在诺特局部环 \(\mathbf{B} = \mathbf{A}_{\mathfrak{p}} / \mathfrak{a}\mathbf{A}_{\mathfrak{p}}\) 中的像。由 § 1, no. 3 的 prop. 7, b)，有

\[
\operatorname{codim} (Z, X) = \dim (B), \quad \operatorname{codim} (Z, X ^ {\prime}) = \dim (B / (\xi_ {1} B + \dots + \xi_ {m} B)).
\]

因此，由 no. 1 的 prop. 2，有

\[
\operatorname{codim} (Z, X ^ {\prime}) \geqslant \operatorname{codim} (Z, X) - m.\tag{11}
\]

若 Z 是 X' 的不可约分支，则 codim(Z, X') = 0，因而 codim(Z, X) ≤ m；这证明了 b)。对 Y 的不可约分支 Z 的集合，在 (11) 的两边分别取下界，即可证明 a)。

反之，假设 \(\operatorname{codim}(\mathbf{Z},\mathbf{X})\leqslant m\)，即 \(\dim (\mathbf{B})\leqslant m\)。由于 \(\mathbf{A}_{\mathfrak{p}}\) 中的每个元素都是 \(\mathbf{A}_{\mathfrak{p}}\) 中的可逆元素与 A 中某个元素的像之积，No. 2 的附论证明存在 A 中的元素 \(x_1^{\prime},\ldots ,x_m^\prime\)，它们在 B 中的像生成 B 的一个定义理想。置 \(\mathrm{H}_i^{\prime} = \mathrm{V}(x_i^{\prime})\)，其中 \(1\leqslant i\leqslant m\)。显然，Z 是 \(\mathbf{X}\cap \mathbf{H}_1^{\prime}\cap \dots \cap \mathbf{H}_m^{\prime}\) 的不可约分支。

推论 1. --- 设 H 是 V 中的非空超曲面。H 在 V 中的余维等于 0 或 1。

codim(H, V) = 1 当且仅当 H 不包含 V 的任何不可约分支。在这种情况下，H 的所有不可约分支在 V 中的余维都是 1。

对于 H 的每个不可约分支 Z，有

\[
0 \leqslant \operatorname{codim} (Z, V) \leqslant 1
\]

根据 prop. 4, b)，并且 codim(Z, V) = 0 当且仅当 Z 是 V 的不可约分支。根据定义，

\[
\operatorname{codim} (\mathrm{H}, \mathrm{V}) = \inf _ {\mathrm{Z}} \operatorname{codim} (\mathrm{Z}, \mathrm{V})
\]

其中 Z 遍历 H 的不可约分支的集合。推论 1 立即由这些备注得到。

推论 2. --- 设 X 是 V 的闭不可约子集，H 是 V 中的超曲面。只有以下三种情形：

\begin{enumerate}
\def\labelenumi{\arabic{enumi})}
\item
  有 \(\mathbf{X}\subset \mathbf{H}\)
\item
  集合 \(\mathbf{X} \cap \mathbf{H}\) 非空，且它的每个不可约分支 \(Z\) 都满足 \(\operatorname{codim}(\mathbf{Z}, \mathbf{X}) = 1\)；
\item
  集合 X ∩ H 为空。
\end{enumerate}

假设 \(X' = X \cap H\) 非空且不同于 \(X\)；每个不可约分支 \(Z\) 属于 \(X'\)，都不同于 \(X\)，且满足 codim \((Z, X) \leqslant 1\)，由 prop. 4, b)。推论 2 立即由此得到。

推论 3. --- 若 A 是唯一分解环（VII, § 3, no. 3, prop. 2），则 A 的高度为 1 的素理想正是由 A 的极元生成的主理想。若 A 还是局部环，则对于 A 的每个高度为 1 的素理想 p，都有 dim(A/p) = dim(A) - 1。

设 \(x\) 是 A 的极元。则 \(Ax\) 是素理想，因为 \(x\) 是极元；它的高度为 1，因为 A 是整环（no. 1, prop. 1）。设 \(p\) 是 A 的高度为 1 的素理想。则 \(V(p)\) 是超曲面 \(V(Ax)\) 的一个不可约分支，对于适当的 \(x\)（prop. 4, c)）。设 \(x = \prod_{i} y_{i}^{n_i}\) 是把 \(x\) 分解为极元之积的分解，

其中 \(y_{i}\) 与 \(y_{j}\) 当 \(i \neq j\) 时互素。\(V(Ax)\) 的不可约分支是 \(V(Ay_{i})\)。因此对于适当的 i，有 \(p = Ay_{i}\)。最后一个断言由 no. 1 的 prop. 2 的 cor. 2 得到。

备注。--- 1) 假设 A 是维数为 d 的局部诺特整环；令 x 为 \(m_{A}\) 中的非零元素，并令 \(H = V(x)\)。由 prop. 4 的 cor. 2，H 的每个不可约分支在 X 中的余维为 1，因而维数 \(\leqslant d - 1\)（§ 1, no. 2, prop. 3）。由 no. 1 的 prop. 2 的 cor. 2，H 的维数为 d - 1，因而这些分支中有一个维数为 d - 1；若 A 是串链环，则它们全都如此。然而，一般而言，H 可能有一个维数 \textless{} d - 1 的不可约分支（cf. p. 87, exercise 7）。

\begin{enumerate}
\def\labelenumi{\arabic{enumi})}
\setcounter{enumi}{1}
\item
  设 \(x_1, \ldots, x_n\) 是 A 中的元素，并置 \(\mathrm{H}_i = \mathrm{V}(x_i)\)，其中 \(1 \leqslant i \leqslant n\)。假设存在一个支集为 \(\mathrm{V} = \operatorname{Spec}(\mathrm{A})\) 的有限生成 A-模 M，使 \((x_1, \ldots, x_n)\) 是 M 的完全截切序列。则 \(\mathrm{H}_1 \cap \ldots \cap \mathrm{H}_n\) 的每个不可约分支在 V 中的余维都是 \(n\)：这由 No. 2 的 Prop. 3 的推论容易得到。
\item
  若诺特环 A 的理想 a 由 m 个元素生成，则 ht(a) ≤ m。这立即由 prop. 4 得到。
\end{enumerate}

命题 5. --- 设 A 为诺特环，\(\mathfrak{p} \subset \mathfrak{q}\) 是 A 的一个非饱和素理想链。由 A 的素理想 \(\mathfrak{r}\) 中满足 \(\mathfrak{p} \subset \mathfrak{r} \subset \mathfrak{q}\) 者组成的集合 E 是无限链。我们有 \(\bigcup \mathfrak{r} = \mathfrak{q}\) 且 \(\bigcap \mathfrak{r} = \mathfrak{p}\)。

将 A 替换为 A/p 后，可归结为 p = \{0\} 的情形。

由 Section 1 的引理 1，有 q = \(\bigcup_{r \in E} r\)，而 II, § 1, no. 1 的 prop. 2 表明 E 是无限的。

令 \(y \neq 0\) 为 \(\bigcap_{r \in E} r\) 中的元素。q 的高度是有限的（\(n^{\circ}1\)，prop. 2 的 cor. 1），并且根据假设有 \(ht(q) \geqslant 2\)。因此存在一个高度为 2 的素理想 \(q' \subset q\)。将证明的第一部分应用于 \(q'\)，可知集合 \(E'\)（由包含于 \(q'\) 的高度为 1 的素理想组成）是无限的；根据假设，这些理想中的每一个都包含 y。现在，根据 n° 1 的 prop. 2 的 cor. 2，诺特局部环 \(B = A_{q'} / yA_{q'}\) 的维数为 1。因此，对每个 \(r \in E'\)，B 的素理想 \(r/yA_{q'}\) 是极小的；于是诺特环 B 有无穷多个极小素理想，这是矛盾的（II, § 4, n° 3, cor. 3 of prop. 14）。所以有 \(\bigcap_{r \in E} r = \{0\}\)。

命题 6.--- 设 A 是维数 \(\geqslant 2\) 的诺特环，\(h\) 是满足 \(0 < h < \dim(A)\) 的整数。

\begin{enumerate}
\def\labelenumi{\alph{enumi})}
\item
  A 有无穷多个高度为 h 的素理想。
\item
  若 A 的维数有限，则 A 有无穷多个素理想 p，使 \(\operatorname{ht}(\mathfrak{p}) = h\) 且 \(\dim(\mathbf{A}/\mathfrak{p}) = \dim(\mathbf{A}) - h\)。
\end{enumerate}

由于 A 的维数是 A 的素理想的高度的上界（§ 1, n° 3, prop. 8），由于 A 的每个素理想都有有限高度（n° 1, prop. 2 的 cor. 1），又由于 \(h < \dim(A)\)，存在整数 \(n > h\) 以及素理想 \(p\)（高度为 \(n\)），因而存在素理想链 \(p_0 \subset \ldots \subset p_n = p\)，其长度为 \(n\)。有 \(ht(p_i) = i\)，其中 \(0 \leqslant i \leqslant n\)，所以有 \(ht(r) = h\)，对于 A 中每个素理想 \(r\)，满足 \(p_{h-1} \subset r \subset p_{h+1}\)。根据 prop. 5，这些理想组成的集合 E 是无限的，故得 \(a\) ).

若 A 的维数有限，则在上述论证中可以设 \(n = \dim(A)\)。对于每个素理想 \(r \in E\)，有 \(\operatorname{ht}(r) = h\)，从而 \(\dim(A/r) \leqslant n - h\)；又由于 \(r \subset p_{h+1} \subset \ldots \subset p_n\) 是长度为 \(n - h\) 的链，有 \(\dim(A/r) = n - h\)。

存在维数为 2 且只有一个高度为 1 的素理想的非诺特整环，例如高度为 2 的赋值环（VI, § 4, n° 4, prop. 5）。

\subsection{环的变换}

命题 7.--- 设 \(\rho: A \to B\) 是诺特局部环的局部同态，M 是有限生成 A-模，N 是有限生成 B-模。置 \(\overline{B} = B \otimes_A \kappa_A = B / \rho(m_A).B\)，以及 \(\overline{N} = N \otimes_B \overline{B} = N / \rho(m_A).N\)。我们有

\[
\dim_ {B} (M \otimes_ {A} N) \leqslant \dim_ {A} (M) + \dim_ {B} (\overline {{N}}),\tag{12}
\]

且当 N 在 A 上平坦时取等号。

不妨设 M 和 N 非零。根据 th. 1 的推论（n° 2），存在 \(m_A\) 的子集 S（分别为 \(m_B\) 的子集 T），使 \(\text{Card}(S) = \dim_A(M)\)（分别有 \(\text{Card}(T) = \dim_{\overline{B}}(\overline{N})\)），并且 M/SM 是有限长度的 A-模（分别地，\(\overline{N}/\overline{T}\overline{N}\) 是有限长度的 B-模）。有 \(\rho(S) \subset m_B\)。令 E 为 B-模 \(M \otimes_A N\)；B-模 \(E/(\rho(S).E + T.E)\) 同构于 M/SM \(\otimes_A N/TN\)，因而具有有限长度；根据 II, § 4, n° 4 的 prop. 18 和 19，有

\[
\begin{array}{r l} \operatorname{Supp} (\mathrm{M} / \mathrm{SM} \otimes_ {\mathrm{A}} \mathrm{N} / \mathrm{TN}) & = \operatorname{Supp} (\mathrm{N} / \mathrm{TN}) \cap {} ^ {a} \rho^ {- 1} (\operatorname{Supp} (\mathrm{M} / \mathrm{SM})) \\ & = \operatorname{Supp} (\mathrm{N} / \mathrm{TN}) \cap {} ^ {a} \rho^ {- 1} (\operatorname{Supp} (\kappa_ {\mathrm{A}})) \\ & = \operatorname{Supp} (\mathrm{N} / \mathrm{TN} \otimes_ {\mathrm{A}} \kappa_ {\mathrm{A}}) = \operatorname{Supp} (\overline {{\mathrm{N}}} / \mathrm{T} \overline {{\mathrm{N}}}) = \left\{\mathfrak {m} _ {\mathrm{B}} \right\}. \end{array}
\]

根据 th. 1，得到不等式

\[
\dim_ {\mathrm{B}} (\mathrm{E}) \leqslant \operatorname{Card} (\mathrm{S}) + \operatorname{Card} (\mathrm{T}) = \dim_ {\mathrm{A}} (\mathrm{M}) + \dim_ {\overline {{\mathrm{B}}}} (\overline {{\mathrm{N}}}).
\]

现在假设 N 在 A 上平坦，并证明公式 (12) 中取等号。设 a（分别为 b）是 M（分别为 N）的湮灭子。根据 II, § 4, n° 4 的 prop. 18 和 19，有

\[
\begin{array}{r l} \operatorname{Supp} (E) & = \operatorname{Supp} (N) \cap {} ^ {a} \rho^ {- 1} (\operatorname{Supp} (M)) \\ & = V (b) \cap {} ^ {a} \rho^ {- 1} (V (a)) = V (b + \rho (a). B). \end{array}
\]

因此有

\[
\dim_ {B} (M \otimes_ {A} N) = \dim (B / (b + \rho (a). B))\tag{13}
\]

并且

\[
\dim_ {\mathbf {A}} (\mathbf {M}) + \dim_ {\overline {{\mathbf {B}}}} (\overline {{\mathbf {N}}}) = \dim (\mathrm{A} / \mathfrak {a}) + \dim (\mathbf {B} / (\mathfrak {b} + \rho (m _ {\mathbf {A}}). \mathbf {B}))  .\tag{14}
\]

置 A' = A/a，B' = B/(b + ρ(a).B)，并令 ρ':A' → B' 为由 ρ 经过取商得到的局部同态。由于 N 的湮灭子为 b，N 是一个支集等于 Spec(B/b) 的有限生成 B/b-模，并且在 A 上平坦。因此，由 ρ 导出的同态 A → B/b 具有 § 2, n° 1 的性质 (PM)（loc. cit., remark 3）；由标量扩张（loc. cit., prop. 1, a)），可推出 ρ' 具有性质 (PM)。根据 loc. cit. 的 prop. 2，有

\[
\dim (\mathbf {B} ^ {\prime}) \geqslant \dim (\mathbf {A} ^ {\prime}) + \dim (\mathbf {B} ^ {\prime} / \rho^ {\prime} (m _ {\mathbf {A} ^ {\prime}}). \mathbf {B} ^ {\prime})\tag{15}
\]

又由于环 \(\mathbf{B}'/\rho'(\mathfrak{m}_{\mathbf{A}'})\mathbf{B}'\) 同构于 \(\mathbf{B}/(\mathfrak{b} + \rho(\mathfrak{m}_{\mathbf{A}})\mathbf{B})\)，我们的断言由公式 (13)、(14) 和 (15) 得出。

推论 1. --- 设 \(\rho: A \to B\) 是诺特局部环的局部同态。

\begin{enumerate}
\def\labelenumi{\alph{enumi})}
\tightlist
\item
  我们有
\end{enumerate}

\[
\dim (\mathbf {B}) \leqslant \dim (\mathbf {A}) + \dim (\mathbf {B} \otimes_ {\mathbf {A}} \kappa_ {\mathbf {A}})\tag{16}
\]

且当 ρ 使 B 成为平坦的 A-模时取等号。

\begin{enumerate}
\def\labelenumi{\alph{enumi})}
\setcounter{enumi}{1}
\tightlist
\item
  假设 \(\rho\) 使 B 成为平坦的 A-模，并设 S 是 \(\mathfrak{m}_{\mathbf{A}}\) 的子集。则 S 是 A 的截切集，当且仅当 \(\rho(\mathbf{S})\) 是 B 的截切集。
\end{enumerate}

断言 \(a)\) 是 prop. 7 在 \(\mathbf{M} = \mathbf{A}\)、\(\mathbf{N} = \mathbf{B}\) 时的特殊情形。

证明 b)。在上述假设下，有

\[
\dim (\mathbf {B}) = \dim (\mathbf {A}) + \dim (\overline {{\mathbf {B}}})\tag{17}
\]

其中 \(\overline{\mathbf{B}} = \mathbf{B} \otimes_{\mathbf{A}} \kappa_{\mathbf{A}} = \mathrm{B} / \rho(\mathfrak{m}_{\mathbf{A}}) . \mathbf{B}\)。由于 \(\rho\) 是单射（I, § 3, no. 5, prop. 9），有

\[
\operatorname{Card} (\rho (S)) = \operatorname{Card} (S).\tag{18}
\]

最后，B′ = B/ρ(S). B 是 A′ = A/SA 上的平坦模，因而

\[
\dim (\mathrm{B} / \rho (\mathrm{S}). \mathrm{B}) = \dim (\mathrm{A} / \mathrm{SA}) + \dim (\overline {{\mathrm{B}}})\tag{19}
\]

因为 \(\mathbf{B}^{\prime}/\rho(\mathfrak{m}_{\mathbf{A}^{\prime}})\mathbf{B}^{\prime}\) 同构于 \(\overline{B}\)。我们的断言立即由公式 (17)、(18) 和 (19) 得到。

推论 2. — 设 ρ : A → B 是诺特环的同态。则

\[
\dim (\mathbf {B}) \leqslant \dim (\mathbf {A}) + \sup _ {\mathfrak {p} \in \operatorname{Spec} (\mathbf {A})} \dim (\mathbf {B} \otimes_ {\mathbf {A}} \kappa (\mathfrak {p})).
\]

设 q 是 B 的素理想，且 \(\mathfrak{p} = \rho^{-1}(\mathfrak{q})\)。由 cor. 1，有 \(\dim(B_{\mathfrak{q}}) \leqslant \dim(A_{\mathfrak{p}}) + \dim(B_{\mathfrak{q}} \otimes_{A} \kappa(\mathfrak{p}))\)。由此根据 § 1, no. 3 的 prop. 6, b) 得到不等式 \(\dim(B_{\mathfrak{q}}) \leqslant \dim(A) + \dim(B \otimes_{A} \kappa(\mathfrak{p}))\)，再由 § 1, no. 3 的 prop. 8 得到结论。

推论 3. --- 设 A 为诺特环，n 为整数 ≥ 0。则

\[
\dim (\mathrm{A} [ \mathrm{X} _ {1},..., \mathrm{X} _ {n} ]) = \dim (\mathrm{A}) + n.
\]

置 \(\mathbf{B} = \mathbf{A}[\mathbf{X}_1, \ldots, \mathbf{X}_n]\)。对于每个素理想 \(\mathfrak{p}\)，即 \(\mathbf{A}\) 的素理想，环 \(\mathbf{B} \otimes_{\mathbf{A}} \kappa(\mathfrak{p})\) 是域上的 \(n\) 元多项式环，因而维数为 \(n\)（\(\S 2\)，no 4, cor. 1 to th. 3）。由 cor. 2，有 \(\dim(\mathbf{B}) \leqslant \dim(\mathbf{A}) + n\)；反向不等式由 \(\S 1\) 的 Example 4, no. 3 得到。

推论 4. --- 设 \(\rho: A \to B\) 是诺特环的同态，a 是 A 的理想。我们有不等式

\[
\operatorname{ht} (\rho (\mathfrak {a}). \mathbf {B}) \leqslant \operatorname{ht} (\mathfrak {a})\tag{20}
\]

若映射 \(^a\rho : \text{Spec}(\mathbf{B}) \to \text{Spec}(\mathbf{A})\) 是满射。若 B 是忠实平坦的 A-模，则有 \(\operatorname{ht}(\rho(\mathfrak{a}).\mathbf{B}) = \operatorname{ht}(\mathfrak{a})\)。

若 B 在 A 上忠实平坦，则 \(^{a}\rho\) 是满射（II, § 2, n° 5, cor. 4 to prop. 11），并且有 ht(a) ≤ ht(ρ(a).B)（§ 2, n° 1, corollary of prop. 2）。

因此，只需在 \(^a\rho\) 为满射的假设下证明不等式 (20)。设 \(\mathfrak{p}\) 是 A 的素理想，满足 \(\mathfrak{a} \subset \mathfrak{p}\) 且 \(\mathrm{ht}(\mathfrak{a}) = \mathrm{ht}(\mathfrak{p})\)（\(\S 1, n^o 3, \text{prop. } 7\)）。令 \(\overline{\mathbf{B}} = \mathbf{B} \otimes_{\mathbf{A}} \kappa(\mathfrak{p})\)，并以 \(h\) 表示从 B 到 \(\overline{\mathbf{B}}\) 的典范同态。若 \(X = ^a\rho^{-1}(\mathfrak{p})\) 是位于 \(\mathfrak{p}\) 之上的 B 的素理想组成的非空集合，我们知道（\(\S 2, n^o 1, \text{lemme 1}\)），映射 \(^ah\) 是从 \(\operatorname{Spec}(\overline{\mathbf{B}})\) 到 \(\operatorname{Spec}(\mathbf{B})\) 的子空间 X 的同胚。设 \(\mathfrak{q}\) 是 \(^ah\) 下来自 \(\overline{\mathbf{B}}\) 的一个极小素理想的像；有 \(\dim(\mathbf{B}_{\mathfrak{q}} \otimes_{\mathbf{A}} \kappa(\mathfrak{p})) = 0\)，并且 Cor. 1 蕴含不等式 \(\dim(\mathbf{B}_{\mathfrak{q}}) \leqslant \dim(\mathbf{A}_{\mathfrak{p}})\)。最后，

\[
\operatorname{ht} (\rho (a). B) \leqslant \operatorname{ht} (q) = \dim (B _ {q}) \leqslant \dim (A _ {p}) = \operatorname{ht} (p) = \operatorname{ht} (a),
\]

从而得到该推论。

命题 8.--- 设 A 为诺特环，a 为 A 的理想，M 为有限生成 A-模，\(\hat{\mathbf{A}}\) 和 \(\hat{\mathbf{M}}\) 分别是 A 和 M 关于 a-进拓扑的分离完备化。

\begin{enumerate}
\def\labelenumi{\alph{enumi})}
\item
  设 m 是包含 a 的 A 的素理想，且 \(\hat{\mathfrak{m}} = \mathfrak{m}\hat{\mathbf{A}}\)。则 \(\hat{\mathfrak{m}}\) 是 \(\hat{\mathbf{A}}\) 的素理想，并且有 \(\dim_{\hat{\mathbf{A}}_{\hat{\mathfrak{m}}}}(\hat{\mathbf{M}}_{\hat{\mathfrak{m}}}) = \dim_{\mathbf{A}_{\mathfrak{m}}}(\mathbf{M}_{\mathfrak{m}})\)。
\item
  我们有 \(\dim_{\hat{\mathbf{A}}}(\hat{\mathbf{M}}) = \sup_{\mathfrak{m}} \dim_{\mathbf{A}_{\mathfrak{m}}}(\mathbf{M}_{\mathfrak{m}})\)，其中 \(\mathfrak{m}\) 遍历包含 a 的 A 的素理想
\end{enumerate}

（分别地，极大理想）组成的集合。特别地，有 \(\dim_{\hat{\mathbf{A}}}(\hat{\mathbf{M}}) \leqslant \dim_{\mathbf{A}}(\mathbf{M})\)。

\begin{enumerate}
\def\labelenumi{\alph{enumi})}
\item
  由于 \(\hat{\mathbf{A}}/\hat{\mathfrak{m}}\) 可识别为 \(\mathbf{A}/\mathfrak{m}\)，\(\hat{\mathfrak{m}}\) 是 \(\hat{\mathbf{A}}\) 的素理想。由 III, § 3, n° 4 的 th. 3，\(\hat{\mathbf{A}}\) 在 \(\mathbf{A}\) 上平坦，因而 \(\hat{\mathbf{A}}_{\hat{\mathfrak{m}}}\) 在 \(\mathbf{A}_{\mathfrak{m}}\) 上平坦。此外，\(\mathbf{A}\) 到 \(\hat{\mathbf{A}}\) 的典范映射诱导出 \(\mathbf{A}/\mathfrak{a}\) 到 \(\hat{\mathbf{A}}/\mathfrak{a}\hat{\mathbf{A}}\) 的同构，因而也诱导出 \(\mathbf{A}_{\mathfrak{m}}/\mathfrak{mA}_{\mathfrak{m}}\) 到 \(\hat{\mathbf{A}}_{\hat{\mathfrak{m}}}/\mathfrak{m}\hat{\mathbf{A}}_{\hat{\mathfrak{m}}}\) 的同构。将 prop. 7 应用于环 \(\mathbf{A}_{\mathfrak{m}}\) 和 \(\hat{\mathbf{A}}_{\hat{\mathfrak{m}}}\)，以及模 \(\mathbf{M}_{\mathfrak{m}}\) 和 \(\hat{\mathbf{A}}_{\hat{\mathfrak{m}}}\)，即可得到结论；模 \(\mathbf{M}_{\mathfrak{m}} \otimes_{\mathbf{A}_{\mathfrak{m}}} \hat{\mathbf{A}}_{\hat{\mathfrak{m}}}\) 同构于 \(\hat{\mathbf{M}}_{\hat{\mathfrak{m}}}\)（III, loc. cit. and prop. 8）。
\item
  由 III, § 3, n° 4 的 prop. 8，映射 m ↦ m̂ 是从包含 a 的 A 的极大理想集合到 Â 的极大理想集合的双射。断言 b) 由此以及 § 1, n° 4 的 prop. 9 得到。
\end{enumerate}

推论 1. --- 设 A 是 Zariski 环（III, § 3, n° 3, def. 2）。对于每个有限生成 A-模 M，有 \(\dim_{\hat{\mathbf{A}}}(\hat{\mathbf{M}}) = \dim_{\mathbf{A}}(\mathbf{M})\)。

事实上，A 的拓扑是 a-进拓扑，其中 a 是包含于 A 的根基中的理想（loc. cit.），也就是说，a 包含于 A 的每个极大理想中。因此只需应用 prop. 8 的断言 b)。

推论 2. --- 设 A 为诺特环，a 为 A 的理想，\(\hat{A}\) 为 A 关于 a-进拓扑的分离完备化。我们有 dim( \(\hat{A}\) ) \(\leqslant\) dim(A)；当 A 是局部环且 \(a\) 不同于 A 时取等号。

推论 3. --- 设 A 为诺特环，n 为整数 ≥ 0。则

\[
\dim (\mathrm{A} [ [ \mathrm{X} _ {1},..., \mathrm{X} _ {n} ] ]) = \dim (\mathrm{A}) + n.
\]

环 A{[}{[}X₁, \ldots, Xₙ{]}{]} 是多项式环 A{[}X₁, \ldots, Xₙ{]} 关于 \(\mathfrak{a}\)-进拓扑的分离完备化，其中 \(\mathfrak{a}\) 是由 X₁, \ldots, Xₙ 生成的理想；因而有

\[
\dim (\mathrm{A} [ [ \mathrm{X} _ {1},..., \mathrm{X} _ {n} ] ]) \leqslant \dim (\mathrm{A} [ \mathrm{X} _ {1},..., \mathrm{X} _ {n} ])
\]

根据 cor. 2。我们还有

\[
\dim (\mathrm{A} [ \mathrm{X} _ {1},..., \mathrm{X} _ {n} ]) = \dim (\mathrm{A}) + n
\]

根据 prop. 7 的 cor. 3。最后，有

\[
\dim (\mathrm{A}) + n \leqslant \dim (\mathrm{A} [ [ \mathrm{X} _ {1},..., \mathrm{X} _ {n} ] ])
\]

根据 § 1, n° 3 的 example 4。

备注。--- 1) 设 A 为诺特环，a 为 A 的理想。假设 A 关于 a-进拓扑是分离且完备的，并考虑受限形式级数环 R = A \{ X \(_{1}\) , \ldots, X \(_{n}\) \}（III, § 4, n° 2, def. 2）。我们有 dim(R) = dim(A) + n。事实上，R 是环 B = A{[}X \(_{1}\) , \ldots, X \(_{n}\) {]} 关于 aB-进拓扑的完备化，因而 dim(R) ≤ dim(A{[}X \(_{1}\) , \ldots, X \(_{n}\) {]}) = dim(A) + n。反向不等式的证明与形式级数的情形相同。

\begin{enumerate}
\def\labelenumi{\arabic{enumi})}
\setcounter{enumi}{1}
\tightlist
\item
  设 A 是诺特局部环，并将其识别为完备化 \(\hat{\mathbf{A}}\) 的子环；设 B 是包含 A 的 \(\hat{\mathbf{A}}\) 的子环。假设 B 是诺特局部环，且有 \(m_{\mathbf{A}}B = m_{\mathbf{B}}\)。则 A 到 B 的典范嵌入延拓为 \(\hat{\mathbf{A}}\) 到 B 的完备化 \(\hat{\mathbf{B}}\) 的同构（III, § 3, no. 5, prop. 11），因而 dim(B) = dim(A)（prop. 8 的 cor. 2）。特别地，这适用于下述情形。* 设 k 是完备的非离散赋值域，A 是多项式环 \(k[\mathrm{X}_1, ..., \mathrm{X}_n]\) 在由 \(\mathrm{X}_1, ..., \mathrm{X}_n\) 生成的素理想处的局部环，B 是系数在 k 中的 \(\mathrm{X}_1, ..., \mathrm{X}_n\) 的收敛幂级数环。则上述假设得到满足，从而有 dim(B) = n。*
\end{enumerate}

\subsection{截切序列的构造}

命题 9.--- 设 A 为诺特环，M 为有限生成 A-模，a 是关于加法和乘法稳定的 A 的子集，\(\mathfrak{q}_{1},\ldots,\mathfrak{q}_{m}\) 是不包含 a 的 A 的素理想，且 \(r\geqslant1\) 为满足下式的整数：

\[
\operatorname{codim} (\mathrm{V} (\mathfrak {a}) \cap \operatorname{Supp} (\mathrm{M}), \operatorname{Supp} (\mathrm{M})) \geqslant r.\tag{21}
\]

存在一个由 a 中元素组成的序列 \((x_{1}, \ldots, x_{r})\)，它不属于理想 \(q_{1}, \ldots, q_{m}\) 中的任何一个，并且序列 \((x_{1}/1, \ldots, x_{r}/1)\) 由 \(A_{p}\) 中元素组成，是 \(A_{p}\)-模 \(M_{p}\) 的截切序列，对每个满足 \(p \in V(a) \cap \text{Supp}(M)\) 的素理想都成立。

我们对 \(r\) 归纳论证。首先设 \(r = 1\)，并以 \(\Phi\) 表示 \(\operatorname{Supp}(M)\) 的（有限）极小元集合。反设不存在满足陈述条件的元素 \(x_1\)（属于 \(\mathfrak{a}\)）。取 \(x \in \mathfrak{a}\)，且不属于 \(\mathfrak{q}_1 \cup \ldots \cup \mathfrak{q}_m\)。根据假设，存在一个素理想 \(\mathfrak{p}\)，属于 \(V(\mathfrak{a}) \cap \operatorname{Supp}(M)\)，使 \(x/1\)（即 \(x\) 在 \(A_{\mathfrak{p}}\) 中的像）不是 \(M_{\mathfrak{p}}\) 的截切元。令 \(\Psi\) 为由素理想 \(\mathfrak{q} \in \Phi\) 且包含于 \(\mathfrak{p}\) 者组成的集合；则 \(\operatorname{Supp}(M_{\mathfrak{p}})\) 的极小元是素理想 \(\mathfrak{q}A_{\mathfrak{p}}\)，它们属于 \(A_{\mathfrak{p}}\)，其中 \(\mathfrak{q}\) 遍历 \(\Psi\)。由 no. 2 的 prop. 3，存在 \(\mathfrak{q} \in \Psi\)，使 \(x/1 \in \mathfrak{q}A_{\mathfrak{p}}\)，因而 \(x \in \mathfrak{q}\)。换言之，有 \(\mathfrak{a} \subset \mathfrak{q}_1 \cup \ldots \cup \mathfrak{q}_m \cup \bigcup_{\mathfrak{q} \in \Phi} \mathfrak{q}\)。由于有 \(\mathfrak{a} \not\subset \mathfrak{q}_j\) 对 \(1 \leqslant j \leqslant m\) 成立，II, § 1, no. 1 的 prop. 2 证明存在素理想 \(\mathfrak{q} \in \Phi\)，包含 \(\mathfrak{a}\)，因而

\[
\mathrm{V} (\mathfrak {q}) \subset \mathrm{V} (\mathfrak {a}) \cap \operatorname{Supp} (\mathrm{M}).
\]

由于 V(a) 包含 Supp(M) 的一个不可约分支，这与假设矛盾：

\[
\operatorname{codim} (\mathrm{V} (\mathfrak {a}) \cap \operatorname{Supp} (\mathrm{M}), \operatorname{Supp} (\mathrm{M})) \geqslant 1.
\]

现在设 \(r \geqslant 2\)。根据归纳假设，可以找到一个序列 \((x_1, ..., x_{r-1})\)，其元素属于 \(\mathfrak{a}\)，不属于 \(\mathfrak{q}_1 \cup ... \cup \mathfrak{q}_m\)，并且对于每个 \(\mathfrak{p} \in V(\mathfrak{a}) \cap \operatorname{Supp}(M)\)，序列 \((x_1/1, ..., x_{r-1}/1)\) 由 \(A_{\mathfrak{p}}\) 中元素组成，是 \(M_{\mathfrak{p}}\) 的截切序列。置 \(N = M / \sum_{i=1}^{r-1} x_i M\)。只需构造元素 \(x_r\)，属于 \(\mathfrak{a}\)，不属于 \(\mathfrak{q}_1 \cup ... \cup \mathfrak{q}_m\)，并且对于每个 \(\mathfrak{p} \in V(\mathfrak{a}) \cap \operatorname{Supp}(M)\)，其像 \(x_r/1\) 属于 \(A_{\mathfrak{p}}\)，且它是 \(N_{\mathfrak{p}}\) 的截切元。根据证明的第一部分，只需证明以下两个关系：

\begin{enumerate}
\def\labelenumi{(\arabic{enumi})}
\setcounter{enumi}{21}
\tightlist
\item
\end{enumerate}

\[
\mathrm{V} (\mathfrak {a}) \cap \operatorname{Supp} (\mathrm{M}) = \mathrm{V} (\mathfrak {a}) \cap \operatorname{Supp} (\mathrm{N}),\tag{23}
\]

\[
\operatorname{codim} (\mathrm{V} (\mathfrak {a}) \cap \operatorname{Supp} (\mathrm{N}), \operatorname{Supp} (\mathrm{N})) \geqslant 1.
\]

现在有

\[
\operatorname{Supp} (\mathrm{N}) = \mathrm{V} (x _ {1}) \cap \dots \cap \mathrm{V} (x _ {r - 1}) \cap \operatorname{Supp} (\mathrm{M})
\]

根据 II, § 4, no. 4 的 prop. 18 的推论，并且由于 \(x_1\)、\ldots、\(x_{r-1}\) 属于 a，可得 (22)。不等式 (23) 随后由假设 (21) 以及 no. 3 的 prop. 4, a) 得到，其中取

\[
m = r - 1, \quad X = \operatorname{Supp} (M), \quad Y = V (\mathfrak {a}) \cap \operatorname{Supp} (N), \quad H _ {i} = V (x _ {i})
\]

从而 X' = Supp(N)。

推论 1. --- 设 A 为诺特环，M 为有限生成 A-模，\(\mathfrak{p}_1, \ldots, \mathfrak{p}_n\)、\(\mathfrak{q}_1, \ldots, \mathfrak{q}_m\) 为 A 的素理想，且 \(r\) 为整数 \(\geqslant 1\)。假设有 \(\mathfrak{p}_i \notin \mathfrak{q}_j\)，其中 \(1 \leqslant i \leqslant n\) 且 \(1 \leqslant j \leqslant m\)，并且有 \(\dim_{\mathbf{A}_{\mathfrak{p}_i}} (\mathbf{M}_{\mathfrak{p}_i}) \geqslant r\) 对 \(1 \leqslant i \leqslant n\) 成立。则存在一个由 A 中元素组成的序列 \((x_1, \ldots, x_r)\)，它属于所有 \(\mathfrak{p}_i\) 而不属于任何 \(\mathfrak{q}_j\)，并且对于 \(1 \leqslant i \leqslant n\)，\(x_1, \ldots, x_r\) 在 \(\mathbf{A}_{\mathfrak{p}_i}\) 中的像组成 \(M_{\mathfrak{p}_i}\) 的截切序列。

置 \(\mathfrak{a} = \bigcap_{i}\mathfrak{p}_{i}\)。由于 \(\mathbf{M}_{\mathfrak{p}_i}\neq \{0\}\)，有 \(\mathfrak{p}_i\in \mathrm{Supp}(\mathbf{M})\)（\(1\leqslant i\leqslant n\)），从而 \(\mathrm{V}(\mathfrak{a})\subset \mathrm{Supp}(\mathbf{M})\)。我们有

\[
\begin{array}{r l} \operatorname{codim} (\mathrm{V} (\mathfrak {a}) \cap \operatorname{Supp} (\mathrm{M}), \operatorname{Supp} (\mathrm{M})) & = \operatorname{codim} (\mathrm{V} (\mathfrak {a}), \operatorname{Supp} (\mathrm{M})) = \\ & = \inf _ {i} \left(\operatorname{codim} (\mathrm{V} (\mathfrak {p} _ {i}), \operatorname{Supp} (\mathrm{M}))\right) = \inf _ {i} \dim (\mathrm{M} _ {\mathfrak {p} _ {i}}) \geqslant r \end{array}
\]

（§ 1, no. 4, prop. 9），于是应用 prop. 9。

推论 2. --- 设 A 为诺特局部环，M 为非零有限生成 A-模，a 是 \(\mathfrak{m}_{\mathbf{A}}\) 的一个关于加法和乘法稳定的子集，并满足 long(M/aM) \textless{} + ∞。a 中存在一个对于 M 为极大截切的子集。

事实上，有 \(\mathrm{V}(\mathfrak{a}) \cap \operatorname{Supp}(\mathbf{M}) = \operatorname{Supp}(\mathbf{M}/\mathfrak{a}\mathbf{M}) = \{\mathfrak{m}_{\mathbf{A}}\}\)（§ 1, no. 4, remark 1），因而 \(\operatorname{codim}(\mathrm{V}(\mathfrak{a}) \cap \operatorname{Supp}(\mathbf{M}), \operatorname{Supp}(\mathbf{M})) = \dim(\operatorname{Supp}(\mathbf{M})) = \dim_{\mathbf{A}}(\mathbf{M})\)，于是应用 prop. 9。

\section{希尔伯特–塞缪尔级数}

\subsection{环 Z((T))}

设 A 为一个环。赋予 A-模 \(\mathbf{A}^{\mathbf{Z}}\) 离散拓扑的乘积拓扑。所有 \((a_{n})\in\mathbf{A}^{\mathbf{Z}}\)，满足存在 \(n_{0}\in\mathbf{Z}\) 使 \(a_{n}=0\) 对 \(n<n_{0}\) 成立的元素组成 \(\mathbf{A}^{\mathbf{Z}}\) 的子模 B。若 \(a=(a_{n})\in\mathbf{B}\)、\(b=(b_{n})\in\mathbf{B}\)，并置 \(ab=c\)，其中 \(c_{n}=\sum_{i+j=n}a_{i}b_{j}\)，则在 B 上定义 A-代数结构。令 T 为 B 中的元素 \((\theta_{n})\)，满足 \(\theta_{n}=0\) 对 \(n\neq1\) 成立，且 \(\theta_{1}=1\)。则 T 在 B 中可逆；对于 B 的每个元素 \(a=(a_{n})\)，族 \((a_{n}\mathrm{T}^{n})_{n\in\mathbf{Z}}\) 在 \(\mathbf{A}^{\mathbf{Z}}\) 中可求和，并且有

\[
a = \sum_ {n \in \mathbf {Z}} a _ {n} \mathrm{T} ^ {n}.
\]

在本章余下部分，我们用 A((T)) 表示 A-代数 B；它包含形式级数代数 A{[}{[}T{]}{]} 和代数 A{[}T, T \(^{-1}\) {]} 作为子代数；它们的交是多项式代数 A{[}T{]}。

备注。--- 环 A((T)) 自然地可识别为环 A{[}{[}T{]}{]} \(_{T}\) 的分式环，该分式环由环 A{[}{[}T{]}{]} 关于由 T 的幂组成的乘法子集定义。

对于 \(n, p\)，其中二者属于 \(\mathbf{Z}\)，我们按下式定义整数 \(\begin{bmatrix} n \\ p \end{bmatrix}\)：

\[
\left\{ \begin{array}{l} {\left[ \begin{array}{l} n \\ p \end{array} \right] = 0 \quad \text {if} \quad p <   0 \quad \text {or} \quad p > n,} \\ {\left[ \begin{array}{l} n \\ p \end{array} \right] = \binom {n} {p} = \frac {n (n - 1) \dots (n - p + 1)}{p !} \quad \text {if} \quad 0 \leqslant p \leqslant n.} \end{array} \right.\tag{1}
\]

有 \(\begin{bmatrix} n \\ p \end{bmatrix} = \begin{bmatrix} n \\ n - p \end{bmatrix}\)，对于 \(n, p \in \mathbf{Z}\)。

引理 1.--- \(\mathbf{Z}((\mathrm{T}))\) 的元素 1 - T 是可逆的。对于每个整数 \(r > 0\)，有

\[
(1 - \mathrm{T}) ^ {- r} = \sum_ {n \in \mathbf {Z}} \left[ \begin{array}{c} n + r - 1 \\ r - 1 \end{array} \right] \mathrm{T} ^ {n} = \sum_ {n \in \mathbf {N}} \binom {n + r - 1} {r - 1} \mathrm{T} ^ {n}.
\]

事实上，1 - T 在环 \(\mathbf{Z}[[\mathrm{T}]]\) 中可逆，其逆为 \(\sum_{m\geqslant 0}\mathbf{T}^m\)；因此有

\[
(1 - \mathrm{T}) ^ {- r} = (\sum_ {m \geqslant 0} \mathrm{T} ^ {m}) ^ {r} = \sum_ {m _ {1}, \dots , m _ {r} \geqslant 0} \mathrm{T} ^ {m _ {1} + m _ {2} + \dots + m _ {r}},
\]

所述公式由 E, III, p. 44, prop. 15 得出。

令 \(\mathrm{Q}(\mathrm{T}) \in \mathbf{Z}[\mathrm{T}, \mathrm{T}^{-1}]\)，\(r\) 为整数 \(> 0\)，且 \(\mathrm{F} = (1 - \mathrm{T})^{-r}\mathrm{Q} \in \mathbf{Z}((\mathrm{T}))\)。置

\[
\mathrm{Q} (\mathrm{T}) = \sum_ {i \in \mathbf {Z}} a _ {i} \mathrm{T} ^ {i}, \quad \mathrm{F} = \sum_ {n \in \mathbf {Z}} \alpha_ {n} \mathrm{T} ^ {n}.
\]

于是由引理 1，有

\[
\alpha_ {n} = \sum_ {i \in \mathbf {Z}} a _ {i} \left[ \begin{array}{c} n - i + r - 1 \\ r - 1 \end{array} \right] = \sum_ {i \leqslant n} a _ {i} \binom {n - i + r - 1} {r - 1}.
\]

令 \(n_1\) 为 \(\overline{\mathbf{R}}\) 中满足 \(i \in \mathbf{Z}\) 且 \(a_i \neq 0\) 的整数所成集合的上确界。对于每个整数 \(n \geqslant n_1\)，有 \(\alpha_n = \tilde{\alpha}(n)\)，其中 \(\tilde{\alpha}\) 是 \(\mathbf{Q}[\mathbf{X}]\) 中由下式定义的多项式：

\[
\tilde {\alpha} (X) = \frac {1}{(r - 1) !} \sum_ {i \in \mathbf {Z}} a _ {i} \prod_ {j = 1} ^ {r - 1} (X - i + j).
\]

若置 \(c = \mathbf{Q}(1) = \sum_{i \in \mathbf{Z}} a_i\)，则有 \(\tilde{\alpha}(\mathbf{X}) = c\mathbf{X}^{r-1}/(r - 1)! + \theta(\mathbf{X})\)，其中 \(\theta\) 是次数 \(\leqslant r - 2\) 的多项式。因此有

\[
\alpha_ {n} = c \frac {n ^ {r - 1}}{(r - 1) !} + \rho_ {n} n ^ {r - 2},\tag{2}
\]

其中有理数 \(\rho_{n}\) 在 \(n\) 无限增大时趋于一个极限。由此得到关系

\[
\mathrm{Q} (1) = (r - 1)! \lim _ {n \rightarrow \infty} n ^ {1 - r} \alpha_ {n}.\tag{3}
\]

若 \(F = \sum_{n \in \mathbf{Z}} a_n T^n\) 和 \(G = \sum_{n \in \mathbf{Z}} b_n T^n\) 是 \(Z((T))\) 的两个元素，我们用“\(F \leqslant G\)”表示关系“\(a_n \leqslant b_n\) 对每个 \(n \in \mathbf{Z}\) 都成立”。这是一个与 \(Z((T))\) 的环结构相容的序关系（A, VI, p. 18, def. 1）。有 \((1 - T)^{-1} \geqslant 1\)。若 \(Q \in Z[T, T^{-1}]\) 满足 \(\geqslant 0\)，则整数 \(Q(1)\) 为正。

引理 2. --- a) 设 F 是 \(\mathbf{Z}((\mathbf{T}))\) 的非零元素，且存在 \(r \in \mathbf{Z}\)，使 \((1 - \mathrm{T})^r\mathrm{F} \in \mathbf{Z}[\mathrm{T}, \mathrm{T}^{-1}]\)；则 F 可唯一写成 \(\mathrm{F} = (1 - \mathrm{T})^{-d}\) . Q 的形式，其中 \(\mathrm{Q} \in \mathbf{Z}[\mathrm{T}, \mathrm{T}^{-1}]\)、\(\mathrm{Q}(1) \neq 0\)，且 \(d \in \mathbf{Z}\)。若 \(\mathrm{F} \geqslant 0\)，则有 \(\mathrm{Q}(1) > 0\) 且 \(d \geqslant 0\)。

\begin{enumerate}
\def\labelenumi{\alph{enumi})}
\setcounter{enumi}{1}
\tightlist
\item
  设 Q、R 属于 Z{[}T, T \(^{-1}\) {]}，d、d' 属于 Z，且 Q(1) \textgreater{} 0。若
\end{enumerate}

\[
(1 - \mathrm{T}) ^ {- d}. \mathrm{Q} \leqslant (1 - \mathrm{T}) ^ {- d ^ {\prime}}. \mathrm{R},
\]

则或者 \(d < d'\)，或者 \(d = d'\) 且 \(\mathbf{Q}(1) \leqslant \mathbf{R}(1)\)。

\begin{enumerate}
\def\labelenumi{\alph{enumi})}
\item
  可以写成 \(F = (1 - T)^{-r} T^{n} P(T)\)，其中 r、\(n \in Z\)，且 \(P(T) \in Z[T]\)。通过欧几里得除法，可以写成 \(P(T) = (1 - T)^{p} R(T)\)，其中 \(R(T) \in Z[T]\) 且 \(R(1) \neq 0\)。因此 \(F = (1 - T)^{-(r-p)} Q(T)\)，其中 \(Q(T) = T^{n} R(T) \in Z[T, T^{-1}]\) 且 \(Q(1) \neq 0\)。这证明了 d 和 Q 的存在性。此外，若 \((1 - T)^{r} Q(T) = (1 - T)^{s} R(T)\)，其中 r \textgreater{} s 且 Q、R 属于 \(Z[T, T^{-1}]\)，则有 \(R(T) = (1 - T)^{r-s} Q(T)\)，从而 \(R(1) = 0\)；这证明了唯一性。假设 F \(\geqslant 0\)；若 d \textless{} 0，则有 \(F(1) = 0\)，这是不可能的，因为 F 非零且其所有系数均为正；所以 \(d \geqslant 0\)。若 d = 0，则 Q = F \(\geqslant 0\)，从而 Q(1) 为正。若 \(d \geqslant 1\)，则由公式 (3)，Q(1) 为正。这证明了 a)。
\item
  假设 \(d \geqslant d'\)。则 \((1 - T)^{-d}((1 - T)^{d - d'} R - Q) \geqslant 0\)；由于 \(S(T) = (1 - T)^{d - d'} R - Q\) 属于 \(Z[T, T^{-1}]\)，根据前述结果可知 \(S(1) \geqslant 0\)。若 \(d > d'\)，则有 \(S(1) = -Q(1) < 0\)，从而矛盾；若 \(d = d'\)，则有 \(S(1) = R(1) - Q(1)\)，从而 \(Q(1) \leqslant R(1)\)。
\end{enumerate}

\subsection{多项式环上分次模的庞加莱级数}

令 \(H_{0}\) 为一个环，I 为有限集，H 为多项式环 \(\mathrm{H}_{0}\big[(\mathrm{X}_{i})_{i\in\mathbf{I}}\big]\)。对于每个 \(i\in I\)，令 \(d_{i}\) 为整数 \(>0\)。赋予 H 一个 Z 型分次环结构，使 \(H_{0}\) 的元素都是 0 次齐次的，而每个 \(X_{i}\) 都是 \(d_{i}\) 次齐次的。当 \(d_{i}=1\) 对每个 \(i\) 都成立时，就得到多项式环的通常分次。

令 M 为有限生成的分次 H-模，且其所有齐次分量都是有限长度的 \(H_{0}\)-模；称 M 的庞加莱级数为元素 \(P_{M}\)，它属于 \(\mathbf{Z}((T))\) 且满足 \(\mathrm{P}_{\mathrm{M}} = \sum_{n \in \mathbf{Z}} \mathrm{long}_{\mathrm{H}_{0}}(\mathrm{M}_{n}) \cdot \mathrm{T}^{n}\)，并置 \(Q_{M} = P_{M} \cdot \prod_{i \in I} (1 - T^{d_{i}})\)。

定理 1. --- \(\mathbf{Q}_{\mathbf{M}}\) 是 \(\mathbf{Z}((\mathbf{T}))\) 的元素，且属于 \(\mathbf{Z}[\mathbf{T}, \mathbf{T}^{-1}]\)。

将 \(H_{0}\) 除以 M 作为 \(H_{0}\)-模的湮灭子，我们归约到 M 是忠实 \(H_{0}\)-模的情形。若 \(a, b \in Z\) 使得 M 作为 H-模由 \(M' = \sum_{a \leqslant i \leqslant b} M_{i}\) 生成，则 \(M'\) 是非零、有限生成、忠实且具有有限长度的 \(H_{0}\)-模；因此，

环 \(\mathbf{H}_0\) 是阿廷环（A, VIII, § 1, no. 3），因而是诺特环（loc. cit., § 9, no. 1）。所以多项式环 H 也是诺特环（loc. cit., § 1, no. 4）。若 I 为空，则 \(\mathbf{H} = \mathbf{H}_0\)，且整数族 \((\mathrm{long}_{\mathbf{H}_0}(\mathbf{M}_n))_{n \in \mathbf{Z}}\) 的支集有限，因为 M 是有限生成的 \(\mathbf{H}_0\)-模；故定理在此情形成立。现在对集合 I 的基数归纳论证，设 I 非空；取 \(j \in \mathbf{I}\)，并令 \(\mathbf{J} = \mathbf{I} - \{j\}\)。记 H' 为由 \(\mathbf{H}_0\) 以及 \(X_i\)（其中 \(i\) 属于 J）生成的 H 的分次子环；考虑 M 上的倍乘映射 \((\mathbf{X}_j)_\mathbf{M}\)，其比率为 \(X_j\)，以及它的核 R 和余核 S。对于每个 \(n \in \mathbf{Z}\)，有 \(\mathbf{H}_0\)-模的正合序列

\[
0 \rightarrow \mathrm{R} _ {n - d _ {j}} \rightarrow \mathrm{M} _ {n - d _ {j}} \rightarrow \mathrm{M} _ {n} \rightarrow \mathrm{S} _ {n} \rightarrow 0,
\]

因而 \(\mathbf{R}_{n - d_j}\) 和 \(\mathbf{S}_n\) 具有有限长度，并且有

\[
\operatorname{long} _ {\mathrm{H} _ {0}} \left(\mathrm{M} _ {n}\right) - \operatorname{long} _ {\mathrm{H} _ {0}} \left(\mathrm{M} _ {n - d _ {j}}\right) = \operatorname{long} _ {\mathrm{H} _ {0}} \left(\mathrm{S} _ {n}\right) - \operatorname{long} _ {\mathrm{H} _ {0}} \left(\mathrm{R} _ {n - d _ {j}}\right).\tag{4}
\]

由于 M 是诺特环 H 上的有限生成模，H-模 R 和 S 都是有限生成的；由于它们被 \(X_{j}\) 湮灭，所以它们是有限生成的 H'-模。根据归纳假设，\(P_{R}\cdot\prod_{i\in J}(1-T^{d_{i}})\) 和 \(P_{S}\cdot\prod_{i\in J}(1-T^{d_{i}})\) 这两个 \(\mathbf{Z}((T))\) 的元素属于 \(\mathbf{Z}[T,T^{-1}]\)；由 (4) 得到

\[
\mathrm{P} _ {\mathrm{M}} - \mathrm{T} ^ {d _ {j}}. \mathrm{P} _ {\mathrm{M}} = \mathrm{P} _ {\mathrm{S}} - \mathrm{T} ^ {d _ {j}}. \mathrm{P} _ {\mathrm{R}},
\]

即 \((1 - \mathrm{T}^{d_j}).\mathrm{P}_{\mathrm{M}} = \mathrm{P}_{\mathrm{S}} - \mathrm{T}^{d_j}.\mathrm{P}_{\mathrm{R}}\)；于是有

\[
\mathrm{P} _ {\mathrm{M}} \cdot \prod_ {i \in \mathrm{I}} (1 - \mathrm{T} ^ {d _ {i}}) = \mathrm{P} _ {\mathrm{S}} \cdot \prod_ {i \in \mathrm{J}} (1 - \mathrm{T} ^ {d _ {i}}) - \mathrm{T} ^ {d _ {j}} \cdot \mathrm{P} _ {\mathrm{R}} \cdot \prod_ {i \in \mathrm{J}} (1 - \mathrm{T} ^ {d _ {i}}),\tag{5}
\]

从而得到结论。

例 1. --- 假设 \(H_{0}\) 是阿廷环，并取 M = H。用前面的记号，有 R = 0 且 S = H'，所以由 (5) 有 \(Q_{H} = Q_{H'}\)；由于 \(Q_{H_{0}} = \text{long}(H_{0})\)，由归纳可得 \(Q_{H} = \text{long}(H_{0})\)，即

\[
\mathrm{P} _ {\mathrm{H}} = \operatorname{long} (\mathrm{H} _ {0}). \prod_ {i \in \mathrm{I}} (1 - \mathrm{T} ^ {d _ {i}}) ^ {- 1}.\tag{6}
\]

以下设 H 取通常的分次，其中 \(d_{i}=1\) 对每个 \(i\in I\) 都成立，并置 \(r=\operatorname{Card}(\mathbf{I})\)；有 \(\mathrm{P}_{\mathrm{M}}=\mathrm{Q}_{\mathrm{M}}(\mathrm{T}).(1-\mathrm{T})^{-r}\)。置 \(c_{\mathrm{M}}=\mathrm{Q}_{\mathrm{M}}(1)\)。于是由 no. 1 的公式 (2)，有：

推论。--- a) 若 r = 0，则有 long \(_{H_{0}}(M)\) = c \(_{M}\)。

\begin{enumerate}
\def\labelenumi{\alph{enumi})}
\setcounter{enumi}{1}
\item
  若 \(r = 1\)，则有 \(\mathrm{long}_{\mathrm{H}_0}(\mathbf{M}_n) = c_{\mathbf{M}}\)，当 \(n\) 充分大时。
\item
  若 \(r > 1\)，则有 long \(_{\mathrm{H}_0}(\mathbf{M}_n) = c_{\mathbf{M}} \frac{n^{r-1}}{(r - 1)!} + \rho_n n^{r-2}\)，其中 \(\rho_n\) 趋于 \(\mathbf{R}\) 中的一个极限，当 \(n\) 无限增大时。
\end{enumerate}

备注。--- 1) 根据引理 2，整数 \(c_{\mathbf{M}}\) 为正。可能有 \(\mathbf{M} \neq 0\) 而 \(c_{\mathbf{M}} = 0\)（cf. prop. 2）。

\begin{enumerate}
\def\labelenumi{\arabic{enumi})}
\setcounter{enumi}{1}
\tightlist
\item
  设 \(0 \to M' \to M \to M'' \to 0\) 是分次 H-模及 0 次同态的正合序列，且 M 是 H 上有限生成的，并且 \(M_n\) 都是 \(H_0\) 上具有有限长度的模，对于每个 \(n\)。则对于每个 \(n \in \mathbf{Z}\)，有
\end{enumerate}

\[
\operatorname{long} _ {\mathrm{H} _ {0}} \left(\mathrm{M} _ {n}\right) = \operatorname{long} _ {\mathrm{H} _ {0}} \left(\mathrm{M} _ {n} ^ {\prime}\right) + \operatorname{long} _ {\mathrm{H} _ {0}} \left(\mathrm{M} _ {n} ^ {\prime \prime}\right),
\]

因此 \(\mathrm{P_M} = \mathrm{P_{M'}} + \mathrm{P_{M''}}\)、\(\mathrm{Q_M} = \mathrm{Q_{M'}} + \mathrm{Q_{M''}}\)，且 \(c_{\mathbf{M}} = c_{\mathbf{M'}} + c_{\mathbf{M''}}\)。

\begin{enumerate}
\def\labelenumi{\arabic{enumi})}
\setcounter{enumi}{2}
\tightlist
\item
  设 \(\mathbf{M}(p)\) 是由 \(\mathbf{M}\) 在分次中平移 \(p\) 得到的模（A, II, p. 165, example 3）。由于有 \(\mathbf{M}(p)_n = \mathbf{M}_{p+n}\)，所以 \(\mathbf{P}_{\mathbf{M}(p)} = \mathrm{T}^{-p}\mathbf{P}_{\mathbf{M}}\)、\(\mathbf{Q}_{\mathbf{M}(p)} = \mathrm{T}^{-p}\mathbf{Q}_{\mathbf{M}}\)，且 \(c_{\mathbf{M}(p)} = c_{\mathbf{M}}\)。
\end{enumerate}

例。--- 2) 假设 H \(_{0}\) 是阿廷环，且 M 是由 s 个齐次、线性无关的元素生成的分次自由 H-模，这些元素的次数分别为 \(\delta_{1}, ..., \delta_{s}\)。则 M 同构于 H( \(-\delta_{1}\) ) ⊕ \(\cdots\) ⊕ H( \(-\delta_{s}\) )。因此由备注 2、3 和例 1，有

\[
\begin{array}{l} \mathrm {P_ {M}} = \operatorname{long} (\mathrm {H_ {0}}) \left(\sum_ {i = 1} ^ {s} \mathrm{T} ^ {\delta_ {i}}\right) (1 - \mathrm{T}) ^ {- r}, \\ \mathrm {Q_ {M}} = \operatorname{long} (\mathrm {H_ {0}}) \left(\sum_ {i = 1} ^ {s} \mathrm{T} ^ {\delta_ {i}}\right), \\ c _ {\mathrm{M}} = s. \operatorname{long} (\mathrm {H_ {0}}). \end{array}
\]

\begin{enumerate}
\def\labelenumi{\arabic{enumi})}
\setcounter{enumi}{2}
\tightlist
\item
  再次假设 \(H_{0}\) 是阿廷环；令 M 为分次 H-模，并假设存在分次 H-模及 0 次同态的正合序列
\end{enumerate}

\[
0 \to \mathrm{L} _ {n} \to \mathrm{L} _ {n - 1} \to \dots \to \mathrm{L} _ {0} \to \mathrm{M} \to 0,
\]

其中，对于 \(k=0,1,...,n\)，\(\mathbf{L}_{k}\) 是由线性无关的齐次元素生成的分次自由 H-模，这些元素的次数分别为 \(\delta_{k,1},..., \delta_{k,m(k)}\)。于是，由备注 2 和例 2，有

\[
\mathrm{Q} _ {\mathrm{M}} = \operatorname{long} (\mathrm{H} _ {0}). \sum_ {0 \leqslant k \leqslant n} \sum_ {1 \leqslant j \leqslant m (k)} (- 1) ^ {k} \mathrm{T} ^ {\delta_ {k, j}},
\]

\[
c _ {\mathbf {M}} = \operatorname{long} (\mathrm{H} _ {0}). \sum_ {0 \leqslant k \leqslant n} (- 1) ^ {k} m (k).
\]

备注 4. --- 可以证明（p. 88, exerc. 4）：在定理 1 的假设下，\(\mathrm{H}_{0}\)-模 \(\mathrm{Tor}_{j}^{\mathrm{H}}(\mathrm{H}_{0},\mathbf{M})\) 具有有限长度，当 \(j > r\) 时为零，并且有 \(c_{\mathbf{M}} = \sum_{j=0}^{r} (-1)^{j} \mathrm{long}_{\mathrm{H}_{0}}(\mathrm{Tor}_{j}^{\mathrm{H}}(\mathrm{H}_{0},\mathbf{M}))\)。更确切地说，H-模

\(T_{j} = \text{Tor}_{j}^{\text{H}}(\text{H}_{0}, \text{M})\) 自然带有分次结构，并且有

\[
\mathrm{Q} _ {\mathrm{M}} = \sum_ {j = 0} ^ {r} (- 1) ^ {j} \mathrm{P} _ {\mathrm{T} _ {j}}.
\]

命题 1.--- 设 M 是分次 H-模。假设 M 由 \(M_{0}\) 生成，且 \(M_{0}\) 是具有有限长度的 \(H_{0}\)-模。则有

\[
\mathrm{P} _ {\mathrm{M}} \leqslant (1 - \mathrm{T}) ^ {- r} \operatorname{long} _ {\mathrm{H} _ {0}} (\mathrm{M} _ {0}), \quad c _ {\mathrm{M}} \leqslant \operatorname{long} _ {\mathrm{H} _ {0}} (\mathrm{M} _ {0}).
\]

此外，下列条件等价：

\begin{enumerate}
\def\labelenumi{(\roman{enumi})}
\tightlist
\item
  \(c_{\mathbf{M}} = \mathrm{long}_{\mathrm{H_0}}(\mathbf{M}_0)\)；
\end{enumerate}

\[
\text{(ii)} \mathrm{P}_{\mathbf{M}} = \operatorname{long}_{\mathbf{H}_{0}}(\mathbf{M}_{0}) \cdot (1 - \mathrm{T})^{-r},\ \text{即， }\mathbf{M} = \mathbf{M}_{0}\text{ 若 }r=0\text{ 且}
\]

\[
\operatorname{long} _ {\mathrm{H} _ {0}} \left(\mathrm{M} _ {n}\right) = \operatorname{long} _ {\mathrm{H} _ {0}} \left(\mathrm{M} _ {0}\right). \binom {n + r - 1} {r - 1}
\]

对 \(n \in \mathbf{N}\) 成立，当 \(r > 0\) 时；

\begin{enumerate}
\def\labelenumi{(\roman{enumi})}
\setcounter{enumi}{2}
\tightlist
\item
  H-模的典范同态
\end{enumerate}

\[
\varphi : \mathrm{H} \otimes_ {\mathrm{H} _ {0}} \mathrm{M} _ {0} \rightarrow \mathrm{M}
\]

是双射。

令 R 表示 φ 的核。由于 φ 是满射，有

\(P_{M} = P_{H \otimes M_{0}} - P_{R} = \text{long}_{H_{0}}(M_{0})(1 - T)^{-r} - P_{R}\) 且 \(c_{M} = \text{long}_{H_{0}}(M_{0}) - c_{R}\)。条件 (i)、(ii) 和 (iii) 分别等价于 \(c_{R} = 0\)、\(P_{R} = 0\) 和 \(R = 0\)。因此有 (iii) \(\Rightarrow\) (ii) \(\Rightarrow\) (i)，只需证明 \(c_{R} = 0\) 蕴含 \(R = 0\)。假设 \(R \neq 0\)，并令 \(0 = N^{h} \subset N^{h-1} \subset \ldots \subset N^{0} = M_{0}\) 为 \(H_{0}\)-模 \(M_{0}\) 的 Jordan–Hölder 序列。令 \(R^{m}\) 为 \(R\) 与 \(H \otimes_{H_{0}} N^{m}\) 在 \(H \otimes_{H_{0}} M_{0}\) 中的像的交；存在整数 \(m\)，它介于 1 和 \(h\) 之间，使 \(R^{m} \neq R^{m-1}\)。置 \(L = R^{m-1}/R^{m}\)；有 \(0 \leqslant c_{L} \leqslant c_{R}\)，只需证明 \(c_{L} \neq 0\)。现在，若 \(k\) 是 \(H_{0}\) 除以 \(N^{m-1}/N^{m}\) 的极大湮灭理想所得的商域，则 L 可识别为 \(k[(X_{i})_{i\in I}]\) 的非零分次子模。因此 L 包含一个与 \(k[(X_{i})_{i\in I}]\) 的平移模同构的子模；由于 \(c_{k[(X_{i})_{i\in I}]} = 1\)，所以 \(c_{L} \geqslant 1\)（备注 2 和 3），这正是我们要证明的。

根据 A, X, p. 160 的 th. 1，条件 (iii) 意味着 \((\mathbf{X}_{1}, \ldots, \mathbf{X}_{r})\) 是 H-模 M 的完全截切序列。

命题 2. --- 假设 \(H_{0}\) 是域，令 M 为有限生成的分次 H-模。令 K 为 H 的分式域。则 \(c_{M}\) 等于 H-模 M 的秩，即 K-向量空间 \(M \otimes_{H} K\) 的维数。

当 M = H 时这是显然的，因为 \(c_{H} = 1\)。此外，令 \(x \in H\) 为 d 次齐次且非零的元素；有 (H/xH) \(\otimes_{H} K = 0\)；由正合序列

\[
0 \rightarrow \mathrm{H} (- d) \rightarrow \mathrm{H} \rightarrow \mathrm{H} / x \mathrm{H} \rightarrow 0,
\]

以及备注 2 和 3，得到 \(c_{\mathrm{H}/x\mathrm{H}}=0\)。因此，当 M 由单个齐次元素生成时，命题成立。一般情形由此得到，因为每个有限生成分次 H-模都有一个合成序列，其商都是上述形式。

备注 6. -- 在 Prop. 2 的假设下，\(c_{\mathbf{M}} = 0\) 当且仅当 M 是挠 H-模，或者等价地，当且仅当 \(\dim_{\mathrm{H}}(\mathbf{M}) < r\)（\(\S 1, \mathfrak{n}^{\circ} 5\)，Example 4）。

\subsection{良过滤模的希尔伯特–塞缪尔级数}

在本段余下部分，我们使用如下记号：若 \(\mathbf{G} \in \mathbf{Z}((\mathbf{T}))\) 且 \(r \in \mathbf{N}\)，置 \(\mathbf{G}^{(r)} = (1 - \mathbf{T})^{-r}\mathbf{G}\)；特别地，若 \(\mathbf{G} = \sum_{n \in \mathbf{Z}} a_n \mathbf{T}^n\)，则

\[
\mathrm{G} ^ {(1)} = \sum_ {n \in \mathbf {Z}} \left(\sum_ {i \leqslant n} a _ {i}\right) \mathrm{T} ^ {n}.
\]

若 \(\mathbf{G} \geqslant 0\)，则有 \(\mathbf{G}^{(r)} \geqslant 0\) 对每个 \(r \in \mathbf{N}\) 成立。。

设 A 为诺特环，q 为 A 的理想，M 为有限生成 A-模。回忆（III, § 3, no. 1, def. 1），M 上的 q-良过滤是从 Z 到 M 的子模集合的映射 F : n ↦ Fn，满足以下三个条件：

\begin{enumerate}
\def\labelenumi{\alph{enumi})}
\item
  有 \(\mathfrak{q}F_n\subset \mathrm{F}_{n + 1}\subset \mathrm{F}_n\) 对每个 \(n\in \mathbf{Z},\)
\item
  存在 \(n_0 \in \mathbf{Z}\)，使得 \(\mathfrak{q}F_n = \mathrm{F}_{n+1}\) 当 \(n \geqslant n_0\) 时成立；
\item
  存在 \(n_{1} \in \mathbf{Z}\)，使得 \(\mathrm{F}_{n_1} = \mathbf{M}\)。
\end{enumerate}

若 \(n_{0}\) 和 \(n_{1}\) 满足上述条件，则对于所有 \(n\in\mathbf{Z}\)，有

\[
\mathfrak {q} ^ {n - n _ {1}} \mathrm{M} \subset \mathrm{F} _ {n} \subset \mathfrak {q} ^ {n - n _ {0}} \mathrm{M}\tag{7}
\]

（回忆一下，按约定，置 \(\mathfrak{q}^r = \mathbf{A}\)，其中 \(r\leqslant 0\)。）

引理 3.--- 若 F 和 F' 是 M 上的两个 q-良过滤，则存在整数 m，使得 \(F_{n}^{\prime} \subset F_{n-m}\) 对每个 \(n \in Z\) 成立。

事实上，存在 \(n_{2}\)，使得 \(\mathbf{F}_{n}^{\prime}\subset\mathfrak{q}^{n-n_{2}}\mathbf{M}\) 对每个 \(n\) 成立，因而 \(\mathbf{F}_{n}^{\prime}\subset\mathbf{F}_{n-(n_{2}-n_{1})}\) 对每个 \(n\) 成立。

引理 4. --- 设 F 是 M 上的 q-良过滤。若 M/qM 具有有限长度，则对于所有 n ∈ Z，M/F \(_{n+1}\) 和 F \(_{n}\) /F \(_{n+1}\) 都具有有限长度。

用 (7) 中的记号，有 long(M/F \(_{n+1}\) ) ≤ long(M/q \(^{n-n_1+1}\) M)，只需证明对于每个 n，q \(^{n}\) M/q \(^{n+1}\) M 具有有限长度。因此可归结为 q-进过滤的情形。设 (x \(_{1}\) , \ldots, x \(_{r}\) ) 是 A-模 q 的有限生成系，I 是 r 个变量中总次数为 n 的单项式组成的有限集

\(X_{1},...,X_{r}\)。从 \((\mathbf{M}/\mathfrak{q}\mathbf{M})^{\mathbf{I}}\) 到 \(q^{n}M/q^{n+1}M\) 的同态把族 \((u_{m})_{m\in\mathbf{I}}\) 映到元素 \(\sum m(x_{1},...,x_{r})u_{m}\)，并且是满射。由于 M/qM 具有有限长度，\(q^{n}M/q^{n+1}M\) 也具有有限长度。

现在假设 M/qM 具有有限长度。设 F 是 M 上的 q-良过滤。存在 \(n_{1} \in Z\)，使 \(F_{n_{1}} = M\)，从而 \(F_{n} = M\) 对 \(n \leqslant n_{1}\) 成立；因此我们通过置下式定义元素 \(H_{M,F}\) 属于 \(\mathbf{Z}((T))\)：

\[
\mathrm{H} _ {\mathrm{M}, \mathrm{F}} = \sum_ {n \in \mathbf {Z}} \operatorname{long} _ {\mathrm{A} / \mathfrak {q}} (\mathrm{F} _ {n} / \mathrm{F} _ {n + 1}). \mathrm{T} ^ {n} \in \mathbf {Z} ((\mathrm{T})) .\tag{8}
\]

定义 1. --- 我们称 H \(_{M,F}\) 为 A-模 M 的希尔伯特–塞缪尔级数（相对于 q-良过滤 F）。

映射 \(n\mapsto\text{long}_{\mathbf{A}}(\mathbf{F}_{n}/\mathbf{F}_{n+1})\) 通常称为 M 的希尔伯特–塞缪尔函数（相对于 F）。

这尤其适用于 q-进过滤的情形（\(F_{n} = q^{n}M\)）；此时置 \(H_{M,F} = H_{M,q}\)。因此有

\[
\mathrm{H} _ {\mathrm{M}, \mathfrak {q}} = \sum_ {n \in \mathbf {Z}} \operatorname{long} _ {\mathrm{A} / \mathfrak {q}} (\mathfrak {q} ^ {n} \mathrm{M} / \mathfrak {q} ^ {n + 1} \mathrm{M}). \mathrm{T} ^ {n}.\tag{9}
\]

命题 3.--- a) 若 F 是 M 上的 q-良过滤，则有

\[
\mathrm{H} _ {\mathrm{M}, \mathrm{F}} ^ {(1)} = \sum_ {n \in \mathbf {Z}} \operatorname{long} _ {\mathrm{A}} (\mathrm{M} / \mathrm{F} _ {n + 1}). \mathrm{T} ^ {n}.\tag{10}
\]

\begin{enumerate}
\def\labelenumi{\alph{enumi})}
\setcounter{enumi}{1}
\tightlist
\item
  若 F 和 F' 是 M 上的两个 q-良过滤，则存在整数 m，使 \(\mathrm{H}_{\mathbf{M},\mathrm{F}'}^{(1)} \geqslant \mathrm{T}^{m}\mathrm{H}_{\mathbf{M},\mathbf{F}}^{(1)}\)。
\end{enumerate}

部分 \(a)\) 立即由 \(\mathbf{H}_{\mathbf{M},\mathbf{F}}^{(1)}\) 的定义得到；部分 \(b)\) 由 \(a)\) 和引理 3 得到。

定理 2. --- 设 A 为诺特环，q 为 A 的理想，M 为有限生成 A-模，使 M/qM 非零且具有有限长度，F 为 M 上的 q-良过滤。

\begin{enumerate}
\def\labelenumi{\alph{enumi})}
\item
  存在唯一确定的整数 \(d \geqslant 0\) 和元素 \(\mathbf{R}\)，属于 \(\mathbf{Z}[\mathrm{T}, \mathrm{T}^{-1}]\)，使 \(\mathbf{R}(1) > 0\) 且 \(\mathrm{H}_{\mathrm{M,F}} = (1 - \mathrm{T})^{-d}\mathbf{R}\)。
\item
  整数 d 和 R(1) 与所选的 q-良过滤 F 无关。
\item
  考虑分次环 \(\operatorname{gr}(A)\)，其中 \(\operatorname{gr}_{n}(A) = \mathfrak{q}^{n}/\mathfrak{q}^{n+1}\)，以及分次 \(\operatorname{gr}(A)\)-模 \(\operatorname{gr}(M)\)，其中 \(\operatorname{gr}_{n}(M) = F_{n}/F_{n+1}\)。由于有 \(F_{n_{1}} = M\)，且 \(\mathfrak{q}F_{n} = F_{n+1}\) 对 \(n \geqslant n_{0}\) 成立，\(\operatorname{gr}(M)\) 由 \(\bigoplus_{n_{1} \leqslant n \leqslant n_{0}} \operatorname{gr}_{n}(M)\) 生成，因而是有限型的。此外，若 \((x_{1},\ldots,x_{r})\) 是 A-模 \(\mathfrak{q}\) 的有限生成系，则 \(\operatorname{gr}(A)\) 由 \(\operatorname{gr}_{0}(A)\) 以及 \(x_{i}\) 模 \(\mathfrak{q}^{2}\) 的类生成，因而同构于 \(H = (A/\mathfrak{q})[X_{1},\ldots,X_{r}]\) 的一个分次商环。由 no. 2 的 th. 1，有
\end{enumerate}

\[
(1 - \mathrm{T}) ^ {r} \mathrm{H} _ {\mathrm{M}, \mathrm{F}} \in \mathbf {Z} [ \mathrm{T}, \mathrm{T} ^ {- 1} ].
\]

有 \(H_{M,F} \neq 0\)，因此存在唯一确定的 \(d \in N\) 和 \(R \in Z[T, T^{-1}]\)，使 \(\mathrm{R}(1) > 0\) 且 \(\mathrm{H}_{\mathrm{M},\mathrm{F}} = (1 - \mathrm{T})^{-d}. R\)（no. 1 的引理 2）。

\begin{enumerate}
\def\labelenumi{\alph{enumi})}
\setcounter{enumi}{1}
\tightlist
\item
  令 F' 为另一个 q-良过滤，并类似地写成
\end{enumerate}

\[
\mathrm{H} _ {\mathbf {M}, \mathrm{F} ^ {\prime}} = (1 - \mathrm{T}) ^ {- d ^ {\prime}} \mathbf {R} ^ {\prime}.
\]

由 Prop. 3, b)，存在整数 m，使 \((1 - T)^{-d'-1}R' \geqslant T^{m}(1 - T)^{-d-1}R\)。根据 no. 1 的引理 2, b)，这蕴含 \(d' \geqslant d\)，并且若 \(d' = d\)，则 \(R'(1) \geqslant R(1)\)。交换 F 和 \(F'\) 的角色，得到 \(d = d'\) 和 \(R(1) = R'(1)\)。

备注 1. --- 用 \(a\) ) 的记号，写成 \(\mathbf{R} = \sum_{i \in \mathbf{Z}} a_i \mathbf{T}^i\)，并假设 \(d > 0\)。由 no. 1，关系 \(\mathrm{H}_{\mathrm{M},\mathrm{F}} = (1 - \mathrm{T})^{-d}\mathrm{R}\) 也可写成

\[
\operatorname{long} _ {\mathbf {A}} \left(\mathrm{F} _ {n} / \mathrm{F} _ {n + 1}\right) = \sum_ {i \in \mathbf {Z}} a _ {i} \left[ \begin{array}{c} n - i + d - 1 \\ d - 1 \end{array} \right] = \sum_ {i \leqslant n} a _ {i} \binom {n - i + d - 1} {d - 1}.\tag{11}
\]

类似地，由于 \(\mathbf{H}_{\mathbf{M},\mathbf{F}}^{(1)} = (1 - \mathbf{T})^{-d - 1}\mathbf{R}\)，有

\[
\operatorname{long} _ {\mathbf {A}} \left(\mathrm{M} / \mathrm{F} _ {n + 1}\right) = \sum_ {i \in \mathbf {Z}} a _ {i} \left[ \begin{array}{c} n - i + d \\ d \end{array} \right] = \sum_ {i \leqslant n} a _ {i} \binom {n - i + d} {d}.\tag{12}
\]

设 A 为诺特环，q 为 A 的理想，M 为有限生成 A-模，且 M/qM 具有有限长度。若 M ≠ qM，则根据定理 2, b)，存在整数 \(d_{q}(M) \geqslant 0\) 和 \(e_{q}(M) > 0\)，使得对于 M 上的每个 q-良过滤，都存在 R ∈ Z{[}T, T \(^{-1}\) {]}，满足

\[
\mathrm{H} _ {\mathbf {M}, \mathrm{F}} = (1 - \mathrm{T}) ^ {- d _ {\mathfrak {q}} (\mathbf {M})} \mathrm{R}, \quad \mathrm{R} (1) = e _ {\mathfrak {q}} (\mathbf {M}).
\]

若 M = qM，则按约定置 \(d_{q}(M) = -\infty\)，\(e_{q}(M) = 0\)。

备注 2. --- 说 M/qM 具有有限长度，意味着

\[
\operatorname{Supp} (\mathrm{M} / \mathfrak {q} \mathrm{M}) = \operatorname{Supp} (\mathrm{M}) \cap \mathrm{V} (\mathfrak {q})
\]

由极大理想组成（IV, § 2, no. 5, prop. 7）。我们将在下文（no. 4, corollary to th. 3）看到，\(d_{\mathfrak{q}}(\mathbf{M})\) 是数 \(\dim_{\mathbf{A}_{\mathfrak{m}}}(M_{\mathfrak{m}})\) 的最小上界，其中 m 遍历集合 \(\operatorname{Supp}(\mathbf{M}) \cap \mathbf{V}(\mathfrak{q})\)。

推论。— 设 A 为诺特环，q 为 A 的理想，M 为有限生成 A-模，使 M/qM 具有有限长度，F 为 M 上的 q-良过滤。

\begin{enumerate}
\def\labelenumi{\alph{enumi})}
\tightlist
\item
  要有 \(d_{\mathfrak{q}}(\mathbf{M}) \leqslant 0\)，必要且充分的条件是序列 \((\mathfrak{q}^{n}\mathbf{M})\) 稳定，或等价地，序列 \((\mathrm{F}_{n})\) 稳定。于是对于所有充分大的 n，
\end{enumerate}

\[
\operatorname{long} \left(\mathrm{M} / \mathrm{F} _ {n + 1}\right) = \operatorname{long} \left(\mathrm{M} / \mathfrak {q} ^ {n + 1} \mathrm{M}\right) = e _ {\mathfrak {q}} (\mathrm{M}).
\]

\begin{enumerate}
\def\labelenumi{\alph{enumi})}
\setcounter{enumi}{1}
\tightlist
\item
  假设 \(d_{\mathfrak{q}}(\mathbf{M}) > 0\)。则有
\end{enumerate}

\begin{enumerate}
\def\labelenumi{(\arabic{enumi})}
\setcounter{enumi}{12}
\tightlist
\item
\end{enumerate}

\[
\operatorname{long} _ {\mathbf {A}} \left(\mathrm{F} _ {n} / \mathrm{F} _ {n + 1}\right) = e _ {\mathfrak {q}} (\mathrm{M}) n ^ {d _ {\mathfrak {q}} (\mathrm{M}) - 1} / \left(d _ {\mathfrak {q}} (\mathrm{M}) - 1\right)! + \rho_ {n} n ^ {d _ {\mathfrak {q}} (\mathrm{M}) - 2},\tag{14}
\]

\[
\operatorname{long} _ {\mathbf {A}} (\mathrm{M} / \mathrm{F} _ {n + 1}) = e _ {\mathfrak {q}} (\mathbf {M}) n ^ {d _ {\mathfrak {q}} (\mathbf {M})} / d _ {\mathfrak {q}} (\mathbf {M})! + \sigma_ {n} n ^ {d _ {\mathfrak {q}} (\mathbf {M}) - 1},
\]

其中 \(\rho_{n}\) 和 \(\sigma_{n}\) 在 \(n\) 无限增大时趋于一个极限。

这立即由 th. 2 和 no. 1 的公式 (2) 得到。

备注 3. --- 假设 q 包含于 A 的根基中。于是根据 Nakayama 引理，序列 (q \(^{n}\) M) 稳定，当且仅当对于充分大的 n 有 q \(^{n}\) M = 0。由推论的 a) 部分随即得到，d \(_{q}\) (M) ≤ 0 当且仅当 M 具有有限长度，并且此时 e \(_{q}\) (M) = long \(_{A}\) (M)。

命题 4. — 设 A 为诺特环，\(x_{1},...,x_{r}\) 是 A 中的元素，\(\mathfrak{x}\) 是它们生成的理想，M 是有限生成 A-模，且 \(M/\mathfrak{x}M\) 非零并具有有限长度。

\begin{enumerate}
\def\labelenumi{\alph{enumi})}
\item
  有 \(d_{\mathfrak{x}}(\mathbf{M}) \leqslant r\)。
\item
  若 \(d_{\mathfrak{x}}(M) = r\)，则 \(e_{\mathfrak{x}}(M) \leqslant \text{long}_{A}(M/\mathfrak{x}M)\)。
\item
  若序列 \((x_{1},...,x_{r})\) 是 M 的完全截切序列（A, X, p. 157），则 \(d_{\mathfrak{x}}(\mathbf{M}) = r\)，且 \(e_{\mathfrak{x}}(\mathbf{M}) = \text{long}_{\mathbf{A}}(\mathbf{M}/\mathfrak{x}\mathbf{M})\)。若 \(x_{i}\) 属于 A 的根基，则逆命题成立。
\end{enumerate}

令 H 为多项式环 \((\mathbf{A} / \mathfrak{x})[\mathbf{X}_1,\dots ,\mathbf{X}_r]\)。我们赋予 \(\mathbf{G} = \bigoplus_{n}\mathfrak{x}^{n}\mathbf{M} / \mathfrak{x}^{n + 1}\mathbf{M}\)

一个分次 H-模结构，其中 \((\mathrm{X}_{i})_{\mathrm{G}}\) 是乘以 \(x_{i}\) 模 \(\mathfrak{x}^{2}\) 的类的映射。采用 No. 2 中的记号 \(P_{G}, Q_{G}, c_{G}\)，有 \(H_{M,\mathfrak{x}} = P_{G}\)，因而 \((1 - T)^{-d_{\mathfrak{x}}(M)}R = (1 - T)^{-r}Q_{G}\)，其中 \(R(1) = e_{\mathfrak{x}}(M) > 0\)，且 \(Q_{G}(1) = c_{G}\)。

因此，要么 \(d_{\mathfrak{x}}(\mathbf{M}) < r\) 且 \(c_{\mathbf{G}} = 0\)，要么 \(d_{\mathfrak{x}}(\mathbf{M}) = r\) 且 \(c_{\mathbf{G}} = e_{\mathfrak{x}}(\mathbf{M})\)。此外，由 no. 2 的 prop. 1，有 \(c_{\mathbf{G}} \leqslant \text{long}(\mathbf{M}/\mathfrak{x}\mathbf{M})\)，且等号成立当且仅当典范同态 \(\mathrm{A}/\mathfrak{x}[\mathrm{X}_{1}, ..., \mathrm{X}_{r}] \otimes_{\mathrm{A}/\mathfrak{x}} \mathrm{M}/\mathfrak{x}\mathrm{M} \to \bigoplus_{n}\mathfrak{x}^{n}\mathrm{M}/\mathfrak{x}^{n+1}\mathrm{M}\) 是双射。

根据 A, X, p. 160 的 th. 1，这就证明了命题。

命题 5. --- 设 \(0 \to \mathbf{M}' \to \mathbf{M} \to \mathbf{M}'' \to 0\) 是诺特环 A 上有限生成模的正合序列，q 是 A 的理想。

\begin{enumerate}
\def\labelenumi{\alph{enumi})}
\item
  若 \(M/\mathfrak{q}M\) 具有有限长度，其充要条件是 \(M'/\mathfrak{q}M'\) 和 \(M''/\mathfrak{q}M''\) 也具有有限长度。
\item
  假设 M/qM 具有有限长度。则属于以下三种情形之一：
\end{enumerate}

\[
1) d _ {\mathfrak {q}} (\mathbf {M}) = d _ {\mathfrak {q}} \left(\mathbf {M} ^ {\prime}\right) > d _ {\mathfrak {q}} \left(\mathbf {M} ^ {\prime \prime}\right) and e _ {\mathfrak {q}} (\mathbf {M}) = e _ {\mathfrak {q}} \left(\mathbf {M} ^ {\prime}\right),
\]

\begin{enumerate}
\def\labelenumi{\arabic{enumi})}
\setcounter{enumi}{1}
\tightlist
\item
  \(d_{\mathfrak{q}}(\mathbf{M}) = d_{\mathfrak{q}}(\mathbf{M}^{\prime \prime}) > d_{\mathfrak{q}}(\mathbf{M}^{\prime})\) 且 \(e_{\mathfrak{q}}(\mathbf{M}) = e_{\mathfrak{q}}(\mathbf{M}^{\prime \prime}),\)
\end{enumerate}

\[
3) d _ {\mathfrak {q}} (\mathbf {M}) = d _ {\mathfrak {q}} \left(\mathbf {M} ^ {\prime}\right) = d _ {\mathfrak {q}} \left(\mathbf {M} ^ {\prime \prime}\right) and e _ {\mathfrak {q}} (\mathbf {M}) = e _ {\mathfrak {q}} \left(\mathbf {M} ^ {\prime}\right) + e _ {\mathfrak {q}} \left(\mathbf {M} ^ {\prime \prime}\right).
\]

\begin{enumerate}
\def\labelenumi{\alph{enumi})}
\item
  有 \(\operatorname{Supp}(\mathbf{M}) = \operatorname{Supp}(\mathbf{M}') \cup \operatorname{Supp}(\mathbf{M}'')\)，断言由备注 2 得到。
\item
  赋予 M 一个 q-良过滤 F（例如 q-进过滤），赋予 M'' 商过滤 F''，赋予 M' 诱导过滤 F'。过滤 F' 和 F'' 都是 q-良过滤（III, § 3, no. 1, prop. 1）。于是对于每个 n，有 A-模的正合序列
\end{enumerate}

\[
0 \rightarrow \mathrm{F} _ {n} ^ {\prime} / \mathrm{F} _ {n + 1} ^ {\prime} \rightarrow \mathrm{F} _ {n} / \mathrm{F} _ {n + 1} \rightarrow \mathrm{F} _ {n} ^ {\prime \prime} / \mathrm{F} _ {n + 1} ^ {\prime \prime} \rightarrow 0
\]

（III, § 2, no. 4, prop. 2），因而有 \(\mathrm{H}_{\mathbf{M},\mathbf{F}} = \mathrm{H}_{\mathbf{M}',\mathbf{F}'} + \mathrm{H}_{\mathbf{M}'',\mathbf{F}''}\)，或等价地

\[
(1 - \mathrm{T}) ^ {- d _ {\mathfrak {q}} (\mathrm{M})} \mathrm{R} = (1 - \mathrm{T}) ^ {- d _ {\mathfrak {q}} (\mathrm{M} ^ {\prime})} \mathrm{R} ^ {\prime} + (1 - \mathrm{T}) ^ {- d _ {\mathfrak {q}} (\mathrm{M} ^ {\prime \prime})} \mathrm{R} ^ {\prime \prime},
\]

其中 \(R,R',R''\in\mathbf{Z}[T,T^{-1}]\)，\(R(1)=e_{\mathfrak{q}}(\mathbf{M})\)，\(R'(1)=e_{\mathfrak{q}}(\mathbf{M}')\)，\(R''(1)=e_{\mathfrak{q}}(\mathbf{M}'')\)。断言 b) 立即由此得到。

\subsection{希尔伯特–塞缪尔函数的次数}

定理 3. --- 设 A 为诺特局部环，q 为不同于 A 的 A 的理想，M 为有限生成 A-模，且 M/qM 具有有限长度。则整数 \(d_{\mathfrak{q}}(\mathbf{M})\) 是 A-模 M 的维数（§ 1, no. 4, def. 8）。

不妨设 \(\mathbf{M} \neq 0\)。先证明不等式 \(d_{\mathfrak{q}}(\mathbf{M}) \leqslant \dim_{\mathbf{A}}(\mathbf{M})\)。由 § 3, no. 5 的 prop. 9 的 cor. 2，存在 \(x_1, ..., x_r \in \mathfrak{q}\)，其中 \(r = \dim_{\mathbf{A}}(\mathbf{M})\)，且 \(\text{long}(\mathbf{M}/\sum_{i=1}^{r} x_i \mathbf{M}) < +\infty\)；置 \(\mathfrak{x} = \sum_{i=1}^{r} x_i \mathbf{A}\)。由 no. 3 的 prop. 4，有 \(d_{\mathfrak{x}}(\mathbf{M}) \leqslant r\)；又有 \(\mathfrak{x} \subset \mathfrak{q}\)，因而 \(\mathrm{H}_{\mathbf{M},\mathfrak{q}}^{(1)} \leqslant \mathrm{H}_{\mathbf{M},\mathfrak{x}}^{(1)}\)，于是（no. 1 的引理 2）

\[
d _ {\mathfrak{q}} (\mathbf {M}) \leqslant d _ {\mathfrak{x}} (\mathbf {M}) \leqslant r = \dim_ {\mathrm{A}} (\mathbf {M}).
\]

现在对 \(\dim_{\mathbf{A}}(\mathbf{M})\) 归纳证明不等式 \(\dim_{\mathbf{A}}(\mathbf{M}) \leqslant d_{\mathfrak{q}}(\mathbf{M})\)；当 \(\dim_{\mathbf{A}}(\mathbf{M}) = 0\) 时显然成立。

假设 \(\dim_{\mathbf{A}}(\mathbf{M}) > 0\)，并且 \(\dim_{\mathbf{A}}(\mathbf{N}) \leqslant d_{\mathfrak{q}}(\mathbf{N})\) 对每个满足 \(\dim_{\mathbf{A}}(\mathbf{N}) < \dim_{\mathbf{A}}(\mathbf{M})\) 的有限生成 A-模 N 都成立。若 \(0 = \mathbf{M}_0 \subset \mathbf{M}_1 \subset \ldots \subset \mathbf{M}_n = \mathbf{M}\) 是 M 的合成序列，则有 \(\dim_{\mathbf{A}}(\mathbf{M}) = \sup(\dim_{\mathbf{A}}(\mathbf{M}_i / \mathbf{M}_{i-1}))\)（§ 1, no. 4, prop. 9），且 \(d_{\mathfrak{q}}(\mathbf{M}) = \sup(d_{\mathfrak{q}}(\mathbf{M}_i / \mathbf{M}_{i-1}))\)（No. 3, prop. 5）。根据 IV, § 1, no. 4 的 th. 1，因此可假设 M 形如 A/p，其中 p 是 A 的素理想；有 \(\mathfrak{p} \ne \mathfrak{m}_{A}\)，因为 \(\dim_{\mathbf{A}}(\mathbf{M}) > 0\)。取 \(x \in \mathfrak{m}_{\mathbf{A}} - \mathfrak{p}\)；M = A/p 上的倍乘映射 \(x_{\mathbf{M}}\) 是单射，并且有正合序列

\[
0 \longrightarrow \mathrm{M} \xrightarrow {x _ {\mathrm{M}}} \mathrm{M} \longrightarrow \mathrm{M} / x \mathrm{M} \longrightarrow 0.
\]

由 § 3, no. 2 的 prop. 3，有 \(\dim_{\mathbf{A}}(\mathbf{M}/x\mathbf{M}) = \dim_{\mathbf{A}}(\mathbf{M}) - 1\)；由 no. 3 的 prop. 5 以及前面的正合序列，有 \(d_{\mathfrak{q}}(\mathbf{M}/x\mathbf{M}) \leqslant d_{\mathfrak{q}}(\mathbf{M}) - 1\)。根据归纳假设，有

\[
\dim_ {\mathbf {A}} (\mathbf {M}) = \dim_ {\mathbf {A}} (\mathbf {M} / x \mathbf {M}) + 1 \leqslant d _ {\mathfrak {q}} (\mathbf {M} / x \mathbf {M}) + 1 \leqslant d _ {\mathfrak {q}} (\mathbf {M}),
\]

这就完成了证明。

推论。--- 设 A 为诺特环，M 为有限生成 A-模，q 是 A 的理想，且 M/qM 具有有限长度。则 \(d_{\mathfrak{q}}(\mathbf{M})\) 是各维数 \(\dim_{\mathrm{A}_{\mathfrak{m}}}\left(\mathrm{M}_{\mathfrak{m}}\right)\) 的上确界，其中 \(\mathfrak{m}\) 遍历有限集 \(S = \operatorname{Supp}(\mathbf{M}) \cap \mathrm{V}(\mathfrak{q})\)；而 \(e_{\mathfrak{q}}(\mathbf{M})\) 是将 \(e_{\mathfrak{q}_{\mathfrak{m}}}\left(\mathrm{M}_{\mathfrak{m}}\right)\) 对满足 \(\mathfrak{m}\in S\) 且 \(\dim_{\mathrm{A}_{\mathfrak{m}}}\left(\mathrm{M}_{\mathfrak{m}}\right) = d_{\mathfrak{q}}(\mathbf{M})\) 的项相加所得的和。

对于每个整数 \(n\)，\(\mathbf{M} / \mathfrak{q}^n\mathbf{M}\) 的长度等于 \(\mathrm{long}_{\mathbf{A}_{\mathfrak{m}}}(\mathbf{M}_{\mathfrak{m}} / \mathfrak{q}_{\mathfrak{m}}^n\mathbf{M}_{\mathfrak{m}})\) 之和（IV, § 2, no. 5, cor. 1 to prop. 7 and corollary to prop. 8）。因此，有 \(\mathbf{H}_{\mathbf{M},\mathfrak{q}} = \sum_{\mathfrak{m}\in S}\mathbf{H}_{\mathbf{M}_{\mathfrak{m}},\mathfrak{q}_{\mathfrak{m}}}\)，从而得到推论。

备注。--- 1) 我们还有 \(d_{\mathfrak{q}}(\mathbf{M}) = \sup_{\mathfrak{m} \in \mathbf{V}(\mathfrak{q})} \dim(\mathbf{M}_{\mathfrak{m}})\)，也就是说，\(d_{\mathfrak{q}}(\mathbf{M}) = \dim(\hat{\mathbf{M}})\)，其中 \(\hat{\mathbf{M}}\) 是 \(\mathbf{M}\) 关于 q-进拓扑的完备化（§ 3, No. 4, prop. 8）。

\begin{enumerate}
\def\labelenumi{\arabic{enumi})}
\setcounter{enumi}{1}
\tightlist
\item
  假设 q 包含于 A 的根基中；则 \(\dim(\tilde{\mathbf{M}})=\dim(\mathbf{M})\)（loc. cit., cor. 1），因而 \(d_{\mathfrak{q}}(\mathbf{M})=\dim(\mathbf{M})\)。
\end{enumerate}

\subsection{商模的希尔伯特–塞缪尔级数}

引理 5. — 设 A 为一个环，M 为 A-模，\((\mathbf{P}_{n})\)、\((\mathbf{Q}_{n})\) 是 M 上由子模组成的两个递降过滤。假设有 \(P_{n} \supset Q_{n}\) 和 \(\text{long}_{\mathbf{A}}(\mathbf{P}_{n}/\mathbf{Q}_{n}) < +\infty\) 对每个 \(n \in Z\) 成立，并且存在整数 \(n_{1}\)，使 \(Q_{n_{1}} = M\)。在 \(\mathbf{Z}((T))\) 中，有不等式

\[
\begin{array}{r l} \sum_ {n \in \mathbf {Z}} \mathrm{long} _ {\mathbf {A}} ((\mathrm{P} _ {n + 1} \cap \mathrm{Q} _ {n}) / \mathrm{Q} _ {n + 1}). & \mathrm{T} ^ {n} \leqslant \sum_ {n \in \mathbf {Z}} \mathrm{long} _ {\mathbf {A}} (\mathrm{P} _ {n + 1} / \mathrm{Q} _ {n + 1}). \\ & \mathrm{T} ^ {n} \leqslant \\ & \leqslant (1 - \mathrm{T}) ^ {- 1} \sum_ {n \in \mathbf {Z}} \mathrm{long} _ {\mathbf {A}} ((\mathrm{P} _ {n + 1} \cap \mathrm{Q} _ {n}) / \mathrm{Q} _ {n + 1}). \end{array}
\]

我们需要证明以下不等式：

\begin{enumerate}
\def\labelenumi{(\arabic{enumi})}
\setcounter{enumi}{14}
\tightlist
\item
\end{enumerate}

\[
\operatorname{long} \left(\left(\mathrm{P} _ {n + 1} \cap \mathrm{Q} _ {n}\right) / \mathrm{Q} _ {n + 1}\right) \leqslant \operatorname{long} \left(\mathrm{P} _ {n + 1} / \mathrm{Q} _ {n + 1}\right),\tag{16}
\]

\[
\operatorname{long} \left(\mathrm{P} _ {n + 1} / \mathrm{Q} _ {n + 1}\right) \leqslant \sum_ {i \leqslant n} \operatorname{long} \left(\left(\mathrm{P} _ {i + 1} \cap \mathrm{Q} _ {i}\right) / \mathrm{Q} _ {i + 1}\right).
\]

第一个不等式是显然的。另一方面，有 \(P_{n+1} \cap Q_i = P_{n+1}\) 对 \(i \leqslant n_1\) 成立，且 \(P_{n+1} \cap Q_{n+1} = Q_{n+1}\)；由此得到不等式

\[
\operatorname{long} \left(\mathrm{P} _ {n + 1} / \mathrm{Q} _ {n + 1}\right) \leqslant \sum_ {i \leqslant n} \operatorname{long} \left(\left(\mathrm{P} _ {n + 1} \cap \mathrm{Q} _ {i}\right) / \left(\mathrm{P} _ {n + 1} \cap \mathrm{Q} _ {i + 1}\right)\right).
\]

但是 A-模 \((\mathrm{P}_{n+1} \cap \mathrm{Q}_i)/(\mathrm{P}_{n+1} \cap \mathrm{Q}_{i+1})\) 同构于 \((\mathrm{P}_{i+1} \cap \mathrm{Q}_i)/\mathrm{Q}_{i+1}\) 的一个子模，因而不等式 (16) 成立。

引理 6. --- 设 A 为一个环，M 为 A-模，且 \((F_{n})\) 是 M 上由子模组成的递降过滤；假设存在整数 \(n_{1}\)，使 \(F_{n_{1}} = M\)。设 f 是 M 的自同态，\(M'\) 是其核，\(M''\) 是其余核。赋予 \(M'\) 过滤 \((F_{n}')\)，它由 \((F_{n})\) 诱导；赋予 \(M''\) 商过滤 \((F_{n}'')\)，它来自 \((F_{n})\)。假设 \(F_{n}/F_{n+1}\) 对每个 \(n \in \mathbf{Z}\) 都具有有限长度，并且存在整数 \(\delta\)，使 \(f(F_{n}) \subset F_{n+\delta}\)。令 \(\varphi\) 为由 f 诱导的 \(\delta\) 次分次自同态，其作用于分次模 \(\operatorname{gr}(M) = \bigoplus_{n} F_{n}/F_{n+1}\)。在 \(\mathbf{Z}((T))\) 的下列元素之间

\[
\mathrm{H} _ {\mathbf {M}} = \sum_ {n \in \mathbf {Z}} \operatorname{long} _ {\mathbf {A}} (\mathrm{F} _ {n} / \mathrm{F} _ {n + 1}). \mathrm{T} ^ {n}
\]

\[
\mathrm{H} _ {\mathbf {M} ^ {\prime}} = \sum_ {n \in \mathbf {Z}} \operatorname{long} _ {\mathbf {A}} (\mathrm{F} _ {n} ^ {\prime} / \mathrm{F} _ {n + 1} ^ {\prime}). \mathrm{T} ^ {n}
\]

\[
\mathrm{H} _ {\mathrm{M} ^ {\prime \prime}} = \sum_ {n \in \mathbf {Z}} \operatorname{long} _ {\mathrm{A}} (\mathrm{F} _ {n} ^ {\prime \prime} / \mathrm{F} _ {n + 1} ^ {\prime \prime}). \mathrm{T} ^ {n}
\]

\[
\mathrm{P} _ {\operatorname{Ker} (\varphi)} = \sum_ {n \in \mathbf {Z}} \operatorname{long} _ {\mathrm{A}} (\operatorname{Ker} (\varphi_ {n})). \mathrm{T} ^ {n},
\]

有不等式

\begin{enumerate}
\def\labelenumi{(\arabic{enumi})}
\setcounter{enumi}{16}
\tightlist
\item
\end{enumerate}

\[
\mathrm{H} _ {\mathbf {M} ^ {\prime}} \leqslant \mathrm{P} _ {\mathbf {K e r} (\varphi)}\tag{18}
\]

\[
(1 - \mathrm{T} ^ {\delta}). \mathrm{H} _ {\mathbf {M}} ^ {(1)} + \mathrm{T} ^ {\delta}. \mathrm{P} _ {\operatorname{Ker} (\varphi)} \leqslant \mathrm{H} _ {\mathbf {M} ^ {\prime \prime}} ^ {(1)} \leqslant (1 - \mathrm{T} ^ {\delta}). \mathrm{H} _ {\mathbf {M}} ^ {(1)} + \mathrm{T} ^ {\delta}. \mathrm{P} _ {\operatorname{Ker} (\varphi)} ^ {(1)}.
\]

M 上的子模序列 \(\mathbf{G}_{n}=f^{-1}(\mathbf{F}_{n+\delta})\) 是递降过滤，并且对于每个整数 n，有 \(F_{n}\subset G_{n}\)。

根据定义，有 \(\operatorname{Ker}(\varphi_n) = (\mathbf{G}_{n+1} \cap \mathbf{F}_n)/\mathbf{F}_{n+1}\)，因而

\[
\mathrm{P} _ {\operatorname{Ker} (\varphi)} = \sum_ {n \in \mathbf {Z}} \operatorname{long} _ {\mathrm{A}} ((\mathrm{G} _ {n + 1} \cap \mathrm{F} _ {n}) / \mathrm{F} _ {n + 1}). \mathrm{T} ^ {n}.\tag{19}
\]

对于每个 \(n\)，A-模 \((\mathbf{M}' \cap \mathbf{F}_n) / (\mathbf{M}' \cap \mathbf{F}_{n+1})\) 可识别为 \((\mathbf{G}_{n+1} \cap \mathbf{F}_n) / \mathbf{F}_{n+1}\) 的一个子模，不等式 (17) 立即由 (19) 得到。此外，根据引理 5，有

\[
\mathrm{P} _ {\operatorname{Ker} (\varphi)} \leqslant \sum_ {n \in \mathbf {Z}} \operatorname{long} _ {\mathrm{A}} \left(\mathrm{G} _ {n + 1} / \mathrm{F} _ {n + 1}\right). \mathrm{T} ^ {n} \leqslant \mathrm{P} _ {\operatorname{Ker} (\varphi)} ^ {(1)}.\tag{20}
\]

对于每个 \(n \in \mathbf{Z}\)，有 A-模的正合序列

\[
0 \longrightarrow \mathrm{G} _ {n + 1} / \mathrm{F} _ {n + 1} \longrightarrow \mathrm{M} / \mathrm{F} _ {n + 1} \xrightarrow {f _ {n}} \mathrm{M} / \mathrm{F} _ {n + \delta + 1} \longrightarrow \mathrm{M} ^ {\prime \prime} / \mathrm{F} _ {n + \delta + 1} ^ {\prime \prime} \longrightarrow 0,
\]

其中 \(f_{n}\) 是由 \(f\) 经过取商得到的。因此有

\[
\operatorname{long} _ {\mathrm{A}} \left(\mathrm{M} ^ {\prime \prime} / \mathrm{F} _ {n + \delta + 1} ^ {\prime \prime}\right) = \operatorname{long} _ {\mathrm{A}} \left(\mathrm{M} / \mathrm{F} _ {n + \delta + 1}\right) - \operatorname{long} _ {\mathrm{A}} \left(\mathrm{M} / \mathrm{F} _ {n + 1}\right) + \operatorname{long} _ {\mathrm{A}} \left(\mathrm{G} _ {n + 1} / \mathrm{F} _ {n + 1}\right).
\]

两边乘以 \(T^{n+\delta}\) 并对 n 求和，得到

\[
\mathrm{H} _ {\mathbf {M} ^ {\prime \prime}} ^ {(1)} = (1 - \mathrm{T} ^ {\delta}) \mathrm{H} _ {\mathbf {M}} ^ {(1)} + \mathrm{T} ^ {\delta}. \sum_ {n \in \mathbf {Z}} \operatorname{long} _ {\mathrm{A}} (\mathrm{G} _ {n + 1} / \mathrm{F} _ {n + 1}). \mathrm{T} ^ {n},\tag{21}
\]

不等式 (18) 立即由 (20) 和 (21) 得到。

引理 7. — 保持引理 6 的记号。

\begin{enumerate}
\def\labelenumi{\alph{enumi})}
\item
  有不等式 \(\mathbf{H}_{\mathbf{M}^{\prime \prime}}^{(1)}\geqslant \frac{1 - \mathrm{T}^{\delta}}{1 - \mathrm{T}}\mathbf{H}_{\mathbf{M}}.\)
\item
  等号成立的充要条件是 φ 为单射。
\item
  在此情况下，有 \(\mathbf{M}^{\prime}\subset \bigcap_{n}\mathbf{F}_{n}\)，且 A-模序列
\end{enumerate}

\[
0 \longrightarrow \operatorname{gr} (\mathrm{M}) \stackrel {\varphi} {\longrightarrow} \operatorname{gr} (\mathrm{M}) \stackrel {v} {\longrightarrow} \operatorname{gr} (\mathrm{M} ^ {\prime \prime}) \longrightarrow 0,
\]

其中 v 是典范映射，是正合的。

断言 a) 和 b) 由引理 6 的公式 (18) 以及定义 \(\mathbf{H}_{\mathbf{M}}^{(1)} = (1 - \mathbf{T})^{-1}.\mathbf{H}_{\mathbf{M}}\) 得到。

假设 \(\varphi\) 是单射。由 III, § 2, no. 8, Th. 1, (i)，有

\[
\operatorname{Ker} (f) \subset f ^ {- 1} (\mathrm{F} _ {n + \delta}) = \mathrm{F} _ {n}
\]

对于所有 n，从而得到 c) 的第一个断言。此外，有正合序列

\[
0 \longrightarrow \mathrm{M} / \mathrm{M} ^ {\prime} \stackrel {f ^ {\prime}} {\longrightarrow} \mathrm{M} \longrightarrow \mathrm{M} / f (\mathrm{M}) \longrightarrow 0,
\]

其中 \(f'\) 是由 \(f\) 经过取商得到的。若 \(\varphi\) 是单射，如上有 \(f'^{-1}(\mathbf{F}_n) = \mathbf{F}_{n-\delta}/\mathbf{M}'\)。在 \(\mathbf{M}/\mathbf{M}'\) 上，由 \(f'\) 的逆像从过滤 \(\mathbf{F}\)（定义在 \(\mathbf{M}\) 上）导出的过滤是 \(n \mapsto \mathbf{F}_{n-\delta}/\mathbf{M}'\)；其伴随分次模为 \(\operatorname{gr}(\mathbf{M})(-\delta)\)，并且有分次模的正合序列（III, § 2, no. 4, prop. 2）

\[
0 \longrightarrow \operatorname{gr} (\mathrm{M}) (- \delta) \xrightarrow {\varphi^ {\prime}} \operatorname{gr} (\mathrm{M}) \longrightarrow \operatorname{gr} \left(\mathrm{M} ^ {\prime \prime}\right) \longrightarrow 0,
\]

其中 \(\varphi_{n}^{\prime} = \varphi_{n - \delta}\) 对每个 \(n\) 都成立。这完成了 \(c\) ) 的证明。

命题 6. — 设 A 为诺特环，M 为有限生成 A-模，q 为 A 的理想，且 M/qM 具有有限长度。设 F 是 M 的 q-良过滤，令 gr(A) = \(\bigoplus_{n\geqslant 0}(\mathfrak{q}^n /\mathfrak{q}^{n + 1})\) 为 A 关于 q-进过滤的伴随分次环。

令 \((x_{1},\ldots,x_{s})\) 为 A 中元素组成的序列，\((\delta_{1},\ldots,\delta_{s})\) 为严格正整数序列，使得 \(x_{i}\in\mathfrak{q}^{\delta_{i}}\) 对 \(1\leqslant i\leqslant s\) 成立，并令 \(\xi_{i}\) 为 \(x_{i}\) 在 \(\mathrm{gr}_{\delta_{i}}(\mathbf{A})=\mathfrak{q}^{\delta_{i}}/\mathfrak{q}^{\delta_{i}+1}\) 中的类。

\begin{enumerate}
\def\labelenumi{\alph{enumi})}
\tightlist
\item
  赋予 A-模 \(\overline{\mathbf{M}} = \mathbf{M}/(x_1\mathbf{M} + \cdots + x_s\mathbf{M})\) F 的商 q-良过滤 \(\overline{\mathbf{F}}\)。于是，在 \(\mathbf{Z}((\mathbf{T}))\) 中有不等式
\end{enumerate}

\[
\mathrm{H} _ {\mathrm{M}, \overline {{\mathrm{F}}}} ^ {(s)} \geqslant \left(\prod_ {i = 1} ^ {s} \frac {1 - \mathrm{T} ^ {\delta_ {i}}}{1 - \mathrm{T}}\right). \mathrm{H} _ {\mathrm{M}, \mathrm{F}}.\tag{22}
\]

\begin{enumerate}
\def\labelenumi{\alph{enumi})}
\setcounter{enumi}{1}
\tightlist
\item
  (22) 中取等号的充要条件是：环 gr(A) 的元素序列 \((\xi_1, \ldots, \xi_s)\) 是模 gr(M) = \(\bigoplus_{n} (\mathbf{F}_n / \mathbf{F}_{n+1})\) 的完全截切序列。
\end{enumerate}

在此情形，从 \(\mathrm{gr}(\mathbf{M}) / \sum_{i=1}^{s} \xi_i \cdot \mathrm{gr}(\mathbf{M})\) 到 \(\mathrm{gr}(\overline{\mathbf{M}}) = \bigoplus_{n} (\overline{\mathbf{F}}_n / \overline{\mathbf{F}}_{n+1})\) 的典范同态是同构。

\begin{enumerate}
\def\labelenumi{\alph{enumi})}
\setcounter{enumi}{2}
\tightlist
\item
  假设 b) 的条件满足，并且每个 A-模 \(\mathbf{M}_{i} = \mathbf{M}/(x_{1}\mathbf{M} + \cdots + x_{i}\mathbf{M})\) 对 \(0 \leqslant i < s\) 都关于 q-进拓扑分离 \footnote{当 \(q\) 包含于 \(A\) 的根基中时尤其如此（III, § 3, no. 3, Prop. 6）。}。则序列 \((x_{1}, ..., x_{s})\) 是 A-模 M 的完全截切序列。
\end{enumerate}

当 \(s = 1\) 时，有 \(\bigcap_{n} F_{n} = \bigcap_{n} q^{n} M\)，且序列 \(\{\xi_{1}\}\) 是 \(\operatorname{gr}(M)\) 的完全截切序列，当且仅当比率为 \(\xi_{1}\) 的倍乘映射在 \(\operatorname{gr}(M)\) 上是单射。将引理 7 应用于倍乘映射 \(f = (x_{1})_{M}\)，它作用于 \(M\)，即可立即得到命题 6。

现在设 \(s \geqslant 2\)，并对 \(s\) 归纳论证。将归纳假设应用于赋予 F 的商过滤 G 的 A-模 \(\mathbf{M}_{1} = \mathbf{M}/x_{1}\mathbf{M}\)，以及序列 \((x_{2}, ..., x_{s})\)，得到不等式

\[
\mathrm{H} _ {\overline {{\mathrm{M,F}}}} ^ {(s - 1)} \geqslant \left(\prod_ {i = 2} ^ {s} \frac {1 - \mathrm{T} ^ {\delta_ {i}}}{1 - \mathrm{T}}\right). \mathrm{H} _ {\mathrm{M} _ {1}, \mathrm{G}};\tag{23}
\]

等号成立当且仅当序列 \((\xi_{2}, \ldots, \xi_{s})\) 是 \(\operatorname{gr}(A)\)-模 \(\operatorname{gr}(\mathbf{M}_{1}) = \bigoplus_{n} \mathrm{G}_{n}/\mathrm{G}_{n+1}\) 的完全截切序列。由于元素 \(\frac{1 - \mathrm{T}^{\delta_{i}}}{1 - \mathrm{T}}\) 属于 \(\mathbf{Z}((\mathrm{T}))\) 且为正，已处理的 \(s = 1\) 情形和公式 (23) 给出不等式

\[
\mathrm{H} _ {\mathrm{M}, \mathrm{F}} ^ {(s)} \geqslant \left(\prod_ {i = 2} ^ {s} \frac {1 - \mathrm{T} ^ {\delta_ {i}}}{1 - \mathrm{T}}\right). \mathrm{H} _ {\mathrm{M} _ {1}, \mathrm{G}} ^ {(1)} \geqslant \left(\prod_ {i = 1} ^ {s} \frac {1 - \mathrm{T} ^ {\delta_ {i}}}{1 - \mathrm{T}}\right). \mathrm{H} _ {\mathrm{M}, \mathrm{F}}.\tag{24}
\]

这证明了 a)。

只有在 (23) 中同时取等号且有等式

\[
\mathrm{H} _ {\mathbf {M} _ {1}, \mathbf {G}} ^ {(1)} = \left(\frac {1 - \mathrm{T} ^ {\delta_ {1}}}{1 - \mathrm{T}}\right) \cdot \mathrm{H} _ {\mathbf {M}, \mathbf {F}}.\tag{25}
\]

后一关系意味着 \(\{\xi_{1}\}\) 是 \(\operatorname{gr}(M)\) 的完全截切序列，并蕴含从 \(\operatorname{gr}(M)/\xi_{1}\operatorname{gr}(M)\) 到 \(\operatorname{gr}(M_{1})\) 的典范同态是同构。换言之，(22) 中取等号，当且仅当 \(\{\xi_{1}\}\) 是 \(\operatorname{gr}(M)\) 的完全截切序列，且 \(\{\xi_{2}, ..., \xi_{s}\}\) 是 \(\operatorname{gr}(M)/\xi_{1}\operatorname{gr}(M)\) 的完全截切序列。这意味着 \(\{\xi_{1}, ..., \xi_{s}\}\) 是 \(\operatorname{gr}(M)\) 的完全截切序列（A, X, p. 160, cor. 2）。因此证明了 \(b\)) 的两个条件等价。假设它们满足；则根据归纳假设，\(\operatorname{gr}(M)\) 可识别为 \(\operatorname{gr}(M_{1})/\sum_{i=2}^{s} \xi_{i}\operatorname{gr}(M_{1})\)；此外，由于 \(\operatorname{gr}(M_{1})\) 可识别为 \(\operatorname{gr}(M)/\xi_{1}\operatorname{gr}(M)\)，\(b\)) 的最后一个断言也得到满足。

现在假设 \(\{\xi_1, \ldots, \xi_s\}\) 是 \(\operatorname{gr}(\mathbf{M})\) 的完全截切序列，且每个 \(\mathbf{M}_i\) 关于 q-进拓扑都是分离的（\(0 \leqslant i < s\)）。由上面的结果和归纳假设，序列 \((x_2, \ldots, x_s)\) 是 \(\mathbf{M}_1\) 的完全截切序列；由于有 \(\mathbf{M}_1 = \mathbf{M}/x_1\mathbf{M}\)，且 \(\{x_1\}\) 是 \(\mathbf{M}\) 的完全截切序列，序列 \((x_1, x_2, \ldots, x_s)\) 是 \(\mathbf{M}\) 的完全截切序列（A, X, p. 160, th. 1）。

\section{正则局部环}
